\documentclass[../main.tex]{subfiles}
\begin{document}

\chapter{GPU Architecture Optimizations - I}
\lessoninfo{
  video={M6, lessons 1--3 (Handling Divergent Branch; Register Optimizations; Unified Register File and Shared Memory)},
  sources={Fung et al.\ \cite{fung2007dwf} (dynamic warp formation); Narasiman et al.\ \cite{narasiman2011twolevel} (large warps); Gebhart et al.\ \cite{gebhart2011energy} (register file cache), \cite{gebhart2011compile} (compiler-managed register hierarchy), \cite{gebhart2012unifying} (unified local memory); Jeon et al.\ \cite{jeon2015rfv} (register file virtualization); NVIDIA A100 \cite{nvidia2020ampere} and H100 \cite{nvidia2022hopper} whitepapers; Harris \cite{harris2013sharedmem} (shared memory latency)},
  textbook={A18 \cite{aamodt2018gpu} ch.~3 (assigned), ch.~4; K12 \cite{kim2012gpgpu} §1.4.1 GPU Pipeline, §1.4.2 Handling Branch Instructions; CPG \cite{nvidia2024cudaguide} §4.1 SIMT Architecture, Compute Capability 7.x--9.0 Shared Memory; PMPP \cite{hwu2022pmpp} ch.~10},
  keywords={divergent branch, divergent warp, reconvergence point, immediate post-dominator, if-conversion, SIMT stack, large warp, dynamic warp formation, SIMD efficiency, live and dead registers, register file cache, register file virtualization, software-managed cache, unified local memory, shared memory carve-out},
}

Lesson~4 followed a warp instruction through the pipeline of an SM and
noted two facts that this lesson builds on: a warp carries an active
mask that can switch lanes off, and its registers sit in a register file
so large that it must be banked.  The first part of this lesson asks
what happens when the threads of one warp disagree about a branch: how
the hardware notices, which lanes run which path, and how the paths are
joined again with a stack of masks and program counters.  Two research
ideas then trade the warp's fixed composition for better lane
utilisation.  The second part looks at the register file itself, why it
is large, slow and power-hungry, and what four observations about
register values suggest: a register file cache in front of it, and a
virtualized file smaller than the names it serves.  The third part
compares registers with shared memory and caches and asks whether one
on-chip pool could serve all three, a research generalisation of the
configurable L1 and shared memory of NVIDIA GPUs.  Like Lesson~7, this
lesson is organised around research papers, and a second aim is to read
such papers.

\section{Handling divergent branches}
\label{sec:l06-divergence}

The SPMD model of Lesson~1 runs one program in every thread, but nothing
forces the threads to follow the same path through it.\source{M6,
divergent branches; K12 §1.4.2 Handling Branch Instructions; A18 ch.~3;
CPG §4.1 SIMT Architecture; \cite{fung2007dwf}.}  This section defines
divergence, shows two ways to execute a divergent branch (predication
and the SIMT stack), and presents two microarchitectural responses to its
cost: large warps and dynamic warp formation.

\subsection{Divergent branches and warps}

The kernel below lets the first two threads of each block do one thing
and the others another; with a warp of four threads, as in the figures
of this lesson, the warp of threads 0--3 splits in half.\margin{The
figures of this lesson draw warps of four threads so that masks stay short; the real
warp size is 32 on every NVIDIA GPU from compute capability 5.0 to 9.0
(CPG Table~21 \cite{nvidia2024cudaguide}).}

\begin{lstlisting}[language=C++, numbers=none]
__global__ void split(float *a) {
  int i = threadIdx.x;      // BB1: compute the condition
  if (i < 2)
    a[i] = a[i] * 2.0f;     // BB2: Work-1
  else
    a[i] = a[i] + 1.0f;     // BB3: Work-2
  a[i] = a[i] - 0.5f;       // BB4: both paths again
}
\end{lstlisting}

The compiler cuts such code into a
\term{basic block}[基本ブロック], a maximal sequence of instructions
that is entered only at its first instruction and left only after its
last, so that either all of its instructions execute or none.  The
blocks and the possible transfers of control between them form the
control-flow graph of \cref{fig:l06-divergent-cfg}, which Lesson~10
defines formally.

\begin{definition}[Divergent branch and divergent warp]
  \label{def:l06-divergent}
  A \term{divergent branch}[ダイバージェント分岐] is a branch whose
  outcome differs among the threads of one warp, so that after it the
  threads of the warp would fetch from different program counters.  The
  warp is then a \term{divergent warp}[ダイバージェントワープ], and the
  phenomenon is called \term{warp divergence}[ワープダイバージェンス].
  A branch on which all threads of every warp agree is warp-uniform and
  costs no more than on a scalar processor.
\end{definition}

\begin{figure}[tb]
  \centering
  % A divergent branch (Lesson 6).  Left: the control-flow graph of
% "if (i < 2) Work-1 else Work-2" for a four-thread warp, each basic
% block with its active mask; the paths split after BB1 and reconverge
% at BB4.  Right: the lanes that execute each block, in time order.
% Redraws the outline figure m6l1_divb (Work-2 -> BB4 edge added).
% Usage: % A divergent branch (Lesson 6).  Left: the control-flow graph of
% "if (i < 2) Work-1 else Work-2" for a four-thread warp, each basic
% block with its active mask; the paths split after BB1 and reconverge
% at BB4.  Right: the lanes that execute each block, in time order.
% Redraws the outline figure m6l1_divb (Work-2 -> BB4 edge added).
% Usage: % A divergent branch (Lesson 6).  Left: the control-flow graph of
% "if (i < 2) Work-1 else Work-2" for a four-thread warp, each basic
% block with its active mask; the paths split after BB1 and reconverge
% at BB4.  Right: the lanes that execute each block, in time order.
% Redraws the outline figure m6l1_divb (Work-2 -> BB4 edge added).
% Usage: \input{figures/l06-divergent-cfg}
% Width: body (120 mm); the drawing is about 108 mm wide.
\begin{tikzpicture}[x=1mm, y=1mm, font=\scriptsize\sffamily,
    >={Stealth[length=1.6mm]},
    bb/.style={draw, minimum width=22mm, minimum height=8mm,
               align=center, inner sep=1pt}]
  % CFG
  \node[bb, fill=ln-accent!20] (b1) at (24,42)
    {BB1: \texttt{if (i < 2)}\\\texttt{1111}};
  \node[bb, fill=ln-source!25] (b2) at (11,26)
    {BB2: Work-1\\\texttt{1100}};
  \node[bb, fill=ln-tip!25] (b3) at (37,26)
    {BB3: Work-2\\\texttt{0011}};
  \node[bb, fill=ln-accent!20] (b4) at (24,10)
    {BB4\\\texttt{1111}};
  \draw[->] (b1) -- node[left, inner sep=2pt] {\iflangja{真}{taken}} (b2);
  \draw[->] (b1) -- node[right, inner sep=2pt] {\iflangja{偽}{not taken}} (b3);
  \draw[->] (b2) -- (b4);
  \draw[->] (b3) -- (b4);
  \node[anchor=north, text=ln-muted, align=center] at (24,4.5)
    {\iflangja{BB4 = 再収束点}{BB4 = reconvergence point}};
  % lanes over time
  \node[anchor=south west, font=\scriptsize\sffamily\bfseries] at (58,47)
    {\iflangja{実行するレーン}{lanes that execute}};
  \foreach \l in {0,...,3}
    \node[text=ln-muted] at (73+\l*7,44) {\iflangja{\l}{lane \l}};
  \foreach \row/\name/\m/\c in {0/BB1/{1,1,1,1}/ln-accent!20, 1/BB2/{1,1,0,0}/ln-source!25,
                                 2/BB3/{0,0,1,1}/ln-tip!25, 3/BB4/{1,1,1,1}/ln-accent!20} {
    \node[anchor=east] at (69,38-\row*9) {\name};
    \foreach \bit [count=\l from 0] in \m {
      \ifnum\bit=1
        \draw[fill=\c] (70.5+\l*7,35.5-\row*9) rectangle ++(5,5);
      \else
        \draw[ln-rule, dashed] (70.5+\l*7,35.5-\row*9) rectangle ++(5,5);
      \fi
    }
  }
  \draw[->, ln-muted] (101,41) -- (101,9) node[midway, right, inner sep=1pt]
    {\iflangja{時間}{time}};
\end{tikzpicture}

% Width: body (120 mm); the drawing is about 108 mm wide.
\begin{tikzpicture}[x=1mm, y=1mm, font=\scriptsize\sffamily,
    >={Stealth[length=1.6mm]},
    bb/.style={draw, minimum width=22mm, minimum height=8mm,
               align=center, inner sep=1pt}]
  % CFG
  \node[bb, fill=ln-accent!20] (b1) at (24,42)
    {BB1: \texttt{if (i < 2)}\\\texttt{1111}};
  \node[bb, fill=ln-source!25] (b2) at (11,26)
    {BB2: Work-1\\\texttt{1100}};
  \node[bb, fill=ln-tip!25] (b3) at (37,26)
    {BB3: Work-2\\\texttt{0011}};
  \node[bb, fill=ln-accent!20] (b4) at (24,10)
    {BB4\\\texttt{1111}};
  \draw[->] (b1) -- node[left, inner sep=2pt] {\iflangja{真}{taken}} (b2);
  \draw[->] (b1) -- node[right, inner sep=2pt] {\iflangja{偽}{not taken}} (b3);
  \draw[->] (b2) -- (b4);
  \draw[->] (b3) -- (b4);
  \node[anchor=north, text=ln-muted, align=center] at (24,4.5)
    {\iflangja{BB4 = 再収束点}{BB4 = reconvergence point}};
  % lanes over time
  \node[anchor=south west, font=\scriptsize\sffamily\bfseries] at (58,47)
    {\iflangja{実行するレーン}{lanes that execute}};
  \foreach \l in {0,...,3}
    \node[text=ln-muted] at (73+\l*7,44) {\iflangja{\l}{lane \l}};
  \foreach \row/\name/\m/\c in {0/BB1/{1,1,1,1}/ln-accent!20, 1/BB2/{1,1,0,0}/ln-source!25,
                                 2/BB3/{0,0,1,1}/ln-tip!25, 3/BB4/{1,1,1,1}/ln-accent!20} {
    \node[anchor=east] at (69,38-\row*9) {\name};
    \foreach \bit [count=\l from 0] in \m {
      \ifnum\bit=1
        \draw[fill=\c] (70.5+\l*7,35.5-\row*9) rectangle ++(5,5);
      \else
        \draw[ln-rule, dashed] (70.5+\l*7,35.5-\row*9) rectangle ++(5,5);
      \fi
    }
  }
  \draw[->, ln-muted] (101,41) -- (101,9) node[midway, right, inner sep=1pt]
    {\iflangja{時間}{time}};
\end{tikzpicture}

% Width: body (120 mm); the drawing is about 108 mm wide.
\begin{tikzpicture}[x=1mm, y=1mm, font=\scriptsize\sffamily,
    >={Stealth[length=1.6mm]},
    bb/.style={draw, minimum width=22mm, minimum height=8mm,
               align=center, inner sep=1pt}]
  % CFG
  \node[bb, fill=ln-accent!20] (b1) at (24,42)
    {BB1: \texttt{if (i < 2)}\\\texttt{1111}};
  \node[bb, fill=ln-source!25] (b2) at (11,26)
    {BB2: Work-1\\\texttt{1100}};
  \node[bb, fill=ln-tip!25] (b3) at (37,26)
    {BB3: Work-2\\\texttt{0011}};
  \node[bb, fill=ln-accent!20] (b4) at (24,10)
    {BB4\\\texttt{1111}};
  \draw[->] (b1) -- node[left, inner sep=2pt] {\iflangja{真}{taken}} (b2);
  \draw[->] (b1) -- node[right, inner sep=2pt] {\iflangja{偽}{not taken}} (b3);
  \draw[->] (b2) -- (b4);
  \draw[->] (b3) -- (b4);
  \node[anchor=north, text=ln-muted, align=center] at (24,4.5)
    {\iflangja{BB4 = 再収束点}{BB4 = reconvergence point}};
  % lanes over time
  \node[anchor=south west, font=\scriptsize\sffamily\bfseries] at (58,47)
    {\iflangja{実行するレーン}{lanes that execute}};
  \foreach \l in {0,...,3}
    \node[text=ln-muted] at (73+\l*7,44) {\iflangja{\l}{lane \l}};
  \foreach \row/\name/\m/\c in {0/BB1/{1,1,1,1}/ln-accent!20, 1/BB2/{1,1,0,0}/ln-source!25,
                                 2/BB3/{0,0,1,1}/ln-tip!25, 3/BB4/{1,1,1,1}/ln-accent!20} {
    \node[anchor=east] at (69,38-\row*9) {\name};
    \foreach \bit [count=\l from 0] in \m {
      \ifnum\bit=1
        \draw[fill=\c] (70.5+\l*7,35.5-\row*9) rectangle ++(5,5);
      \else
        \draw[ln-rule, dashed] (70.5+\l*7,35.5-\row*9) rectangle ++(5,5);
      \fi
    }
  }
  \draw[->, ln-muted] (101,41) -- (101,9) node[midway, right, inner sep=1pt]
    {\iflangja{時間}{time}};
\end{tikzpicture}

  \caption{A divergent branch for a warp of four threads.  Left: the
    control-flow graph, each basic block with the active mask under
    which it runs (bit $k$ from the left is lane~$k$).  Right: the
    lanes that execute each block; BB2 and BB3 run one after the other,
    each with half of the lanes idle.}
  \label{fig:l06-divergent-cfg}
\end{figure}

The active mask of Lesson~4 records which lanes take part in each block
(\cref{fig:l06-divergent-cfg}): BB1 runs with \texttt{1111}, BB2 with
\texttt{1100}, BB3 with \texttt{0011} and BB4 again with \texttt{1111}.
The warp has one program counter, so the two paths cannot run at the
same time; they run one after the other, each with the lanes of the
other path switched off, and the two paths' times add up.\caution{A
divergent branch does not make the code wrong, only slower: the result
is the one a scalar machine would compute.  Divergence between
\emph{different} warps costs nothing, since each warp has its own
program counter.}  The point where the paths meet again, BB4, is where
the warp can resume with its full mask.

\begin{definition}[Reconvergence point]
  \label{def:l06-reconvergence}
  The \term{reconvergence point}[再収束点] of a divergent branch is the
  first instruction at which all of the warp's paths from the branch are
  guaranteed to meet again, so that the warp can continue as a whole.
  In compiler terms, a block $Y$ is a
  \term{post-dominator}[後支配節点] of a block $X$ if every path from
  $X$ to the exit of the program passes through $Y$, and the
  \term{immediate post-dominator (IPDOM)}[直接後支配節点]\index{IPDOM|see{immediate post-dominator (IPDOM)}} of $X$ is its
  closest post-dominator other than $X$ itself.  The IPDOM of the
  branch block is the usual reconvergence point; in
  \cref{fig:l06-divergent-cfg} it is BB4.
\end{definition}

\margin{Fung et al.\ report that reconverging at the immediate
post-dominator, rather than letting the diverged parts of a warp run
apart until the end of the kernel, gives a speed-up of 93.4\,\%
\cite{fung2007dwf}.}%
Reconverging as early as possible matters because every block executed
after the branch and before the reconvergence point runs with a partial
mask.

\subsection{If-conversion and predication}

The oldest way to remove a branch is to remove the jump.

\begin{definition}[If-conversion]
  \label{def:l06-if-conversion}
  The transformation called \term{if-conversion}[if変換] replaces a branch and its two paths by
  straight-line code: an instruction computes the condition into a
  \term{predicate}[プレディケート], a one-bit register per lane, and
  each instruction of the former paths is guarded by the predicate or
  its negation (\cref{fig:l06-if-conversion}).  A lane whose guard is
  false does not write the result.  The resulting code is executed by
  the predicated execution of Lesson~5.
\end{definition}

\begin{marginfigure}
  \resizebox{\linewidth}{!}{% If-conversion (Lesson 6): the branch of figure l06-divergent-cfg
% becomes one straight-line block of predicated instructions that every
% lane fetches; no lane ever jumps.
% Redraws the outline figure m6l1_divb2.
% Usage: % If-conversion (Lesson 6): the branch of figure l06-divergent-cfg
% becomes one straight-line block of predicated instructions that every
% lane fetches; no lane ever jumps.
% Redraws the outline figure m6l1_divb2.
% Usage: % If-conversion (Lesson 6): the branch of figure l06-divergent-cfg
% becomes one straight-line block of predicated instructions that every
% lane fetches; no lane ever jumps.
% Redraws the outline figure m6l1_divb2.
% Usage: \input{figures/l06-if-conversion}  (in a marginfigure)
% Width: margin (50 mm); the drawing is about 46 mm wide.
\begin{tikzpicture}[x=1mm, y=1mm, font=\scriptsize\sffamily,
    >={Stealth[length=1.6mm]},
    ins/.style={draw, minimum width=34mm, minimum height=6mm,
                align=center, inner sep=1pt}]
  \node[ins, fill=ln-accent!20] (p) at (0,30)
    {\iflangja{プレディケート\\\texttt{p = (i < 2)}}{set predicate \texttt{p = (i < 2)}}};
  \node[ins, fill=ln-source!25] (a) at (0,21) {\texttt{(p)} Work-1};
  \node[ins, fill=ln-tip!25] (b) at (0,12) {\texttt{(!p)} Work-2};
  \node[ins, fill=ln-accent!20] (c) at (0,3) {BB4};
  \draw[->] (p) -- (a);
  \draw[->] (a) -- (b);
  \draw[->] (b) -- (c);
  \draw[ln-muted] (19,33) -- ++(1.5,0) -- ++(0,-33) -- ++(-1.5,0);
  \node[rotate=90, text=ln-muted] at (23.5,16.5)
    {\iflangja{全レーンがフェッチ}{fetched by every lane}};
\end{tikzpicture}
  (in a marginfigure)
% Width: margin (50 mm); the drawing is about 46 mm wide.
\begin{tikzpicture}[x=1mm, y=1mm, font=\scriptsize\sffamily,
    >={Stealth[length=1.6mm]},
    ins/.style={draw, minimum width=34mm, minimum height=6mm,
                align=center, inner sep=1pt}]
  \node[ins, fill=ln-accent!20] (p) at (0,30)
    {\iflangja{プレディケート\\\texttt{p = (i < 2)}}{set predicate \texttt{p = (i < 2)}}};
  \node[ins, fill=ln-source!25] (a) at (0,21) {\texttt{(p)} Work-1};
  \node[ins, fill=ln-tip!25] (b) at (0,12) {\texttt{(!p)} Work-2};
  \node[ins, fill=ln-accent!20] (c) at (0,3) {BB4};
  \draw[->] (p) -- (a);
  \draw[->] (a) -- (b);
  \draw[->] (b) -- (c);
  \draw[ln-muted] (19,33) -- ++(1.5,0) -- ++(0,-33) -- ++(-1.5,0);
  \node[rotate=90, text=ln-muted] at (23.5,16.5)
    {\iflangja{全レーンがフェッチ}{fetched by every lane}};
\end{tikzpicture}
  (in a marginfigure)
% Width: margin (50 mm); the drawing is about 46 mm wide.
\begin{tikzpicture}[x=1mm, y=1mm, font=\scriptsize\sffamily,
    >={Stealth[length=1.6mm]},
    ins/.style={draw, minimum width=34mm, minimum height=6mm,
                align=center, inner sep=1pt}]
  \node[ins, fill=ln-accent!20] (p) at (0,30)
    {\iflangja{プレディケート\\\texttt{p = (i < 2)}}{set predicate \texttt{p = (i < 2)}}};
  \node[ins, fill=ln-source!25] (a) at (0,21) {\texttt{(p)} Work-1};
  \node[ins, fill=ln-tip!25] (b) at (0,12) {\texttt{(!p)} Work-2};
  \node[ins, fill=ln-accent!20] (c) at (0,3) {BB4};
  \draw[->] (p) -- (a);
  \draw[->] (a) -- (b);
  \draw[->] (b) -- (c);
  \draw[ln-muted] (19,33) -- ++(1.5,0) -- ++(0,-33) -- ++(-1.5,0);
  \node[rotate=90, text=ln-muted] at (23.5,16.5)
    {\iflangja{全レーンがフェッチ}{fetched by every lane}};
\end{tikzpicture}
}
  \caption{The branch of \cref{fig:l06-divergent-cfg} after
    if-conversion: one basic block whose guarded instructions every
    lane fetches.}
  \label{fig:l06-if-conversion}
\end{marginfigure}

If-conversion is standard on vector processors, and CPUs use it for
short branch bodies, where a conditional move is cheaper than a possible
misprediction.  On a GPU it removes divergence altogether: no
lane ever jumps, so the warp keeps one program counter.  The price is
that the instructions of both paths are fetched and issued even when no
lane needs one of them, so the compiler applies if-conversion only to
short bodies.\margin{In PTX the guard is written \texttt{@p} or
\texttt{@!p} in front of the instruction; Lesson~10 shows the
predicate registers.}

\begin{example}
  \label{ex:l06-if-conversion}
  In the kernel above let BB1 and BB4 have 2 instructions each and BB2
  and BB3 have 6 each.  Take a warp of four lanes and count lane-slots,
  one per lane per issued instruction.
  \begin{itemize}
    \item \emph{Divergent warp, branch code}: $2+6+6+2 = 16$ issued
      instructions, $64$ lane-slots, of which $2\cdot4 + 6\cdot2 +
      6\cdot2 + 2\cdot4 = 40$ do work.
    \item \emph{Divergent warp, if-converted}: the same 16 instructions
      (the predicate is computed in BB1) and 40 useful lane-slots.
    \item \emph{Uniform warp (all lanes take BB2), branch code}:
      $2+6+2 = 10$ instructions, all lanes busy.
    \item \emph{Uniform warp, if-converted}: still 16 instructions; the
      6 of BB3 do nothing in any lane.
  \end{itemize}
  If-conversion makes the cost independent of the data.  In this
  count it ties when the warp diverges and loses when the warp is
  uniform; what it saves is the overhead the count leaves out (the
  branch instruction, the stack push and pop, the reconvergence check
  and the fetch redirect), which matters only for short bodies.\tip{Cost of
if-conversion: the sum of both bodies, always.  Cost of a branch: the
body taken if the warp is uniform, the sum of both if it diverges.}
\end{example}

\subsection{The SIMT stack}
\label{sec:l06-simt-stack}

Without predication the hardware itself must run both paths and join
them.\caution{Lecture vs textbook: the two-field entry built below leaves implicit
how the full mask returns at the reconvergence point.  K12 §1.4.2 gives
the entry three fields, \enquote{the PC address of the reconvergence
point, active mask bits, and the next PC address}
\cite{kim2012gpgpu}.}  The mechanism is easiest to understand by building it in steps,
following the warp of \cref{fig:l06-divergent-cfg}.
\begin{enumerate}
  \item After BB1 the next PCs of the four threads differ (BB2 for
    lanes 0--1, BB3 for lanes 2--3): the branch is divergent.  One PC
    and one mask cannot express both destinations.
  \item Following one path only (BB1, BB2, BB4) is wrong: the threads
    of lanes 2--3 never execute BB3.
  \item The fix is to remember the path not taken: an \emph{alternate
    PC} (BB3) with its mask (\texttt{0011}), executed after the taken
    path.  The order becomes BB1, BB2, BB3, BB4.
  \item One alternate slot is not enough when branches nest.  If BB2
    itself diverges (\cref{fig:l06-simt-stack}, left: BB4 with
    \texttt{1000}, BB5 with \texttt{0100}), its alternate PC BB5 would
    overwrite BB3 in the single slot, and BB3 would be lost.
  \item The pending paths therefore go onto a stack.
\end{enumerate}

\begin{definition}[SIMT stack]
  \label{def:l06-simt-stack}
  The \term{SIMT stack}[SIMTスタック] of a warp holds one entry per
  pending path of a divergent branch: the alternate PC with its active
  mask, and the reconvergence PC at which the current path ends.  On a
  divergent branch the hardware pushes an entry for the path not taken
  and continues with the taken path under the reduced mask.  Before
  every fetch it compares the next PC with the reconvergence PC of the
  top entry; if they are equal, it \emph{pops} the entry and fetches
  from the alternate PC with the entry's mask instead.  A
  \term{nested branch}[入れ子の分岐], a divergent branch inside a path
  of another, simply pushes another entry.
\end{definition}

\begin{figure}[tb]
  \centering
  % The SIMT stack on a nested branch (Lesson 6).  Left: a CFG in which
% BB2 branches again (BB4 1000, BB5 0100) and all paths reconverge at
% BB6.  Right: the stack of (alternate PC with its mask, reconvergence
% PC) entries after BB1, after BB2, after the pop at the end of BB4 and
% after the pop at the end of BB5 (top entry first).
% Redraws the outline figures m6l1_gpuflow4 and m6l1_gpuflow5 (BB2 mask
% corrected to 1100).
% Usage: % The SIMT stack on a nested branch (Lesson 6).  Left: a CFG in which
% BB2 branches again (BB4 1000, BB5 0100) and all paths reconverge at
% BB6.  Right: the stack of (alternate PC with its mask, reconvergence
% PC) entries after BB1, after BB2, after the pop at the end of BB4 and
% after the pop at the end of BB5 (top entry first).
% Redraws the outline figures m6l1_gpuflow4 and m6l1_gpuflow5 (BB2 mask
% corrected to 1100).
% Usage: % The SIMT stack on a nested branch (Lesson 6).  Left: a CFG in which
% BB2 branches again (BB4 1000, BB5 0100) and all paths reconverge at
% BB6.  Right: the stack of (alternate PC with its mask, reconvergence
% PC) entries after BB1, after BB2, after the pop at the end of BB4 and
% after the pop at the end of BB5 (top entry first).
% Redraws the outline figures m6l1_gpuflow4 and m6l1_gpuflow5 (BB2 mask
% corrected to 1100).
% Usage: \input{figures/l06-simt-stack}
% Width: body (120 mm); the drawing is about 112 mm wide.
\begin{tikzpicture}[x=1mm, y=1mm, font=\scriptsize\sffamily,
    >={Stealth[length=1.6mm]},
    bb/.style={draw, minimum width=13mm, minimum height=7.5mm,
               align=center, inner sep=1pt}]
  % nested CFG
  \node[bb, fill=ln-accent!20] (b1) at (22,56) {BB1\\\texttt{1111}};
  \node[bb, fill=ln-source!25] (b2) at (12,41) {BB2\\\texttt{1100}};
  \node[bb, fill=ln-tip!25] (b3) at (36,41) {BB3\\\texttt{0011}};
  \node[bb, fill=ln-source!12] (b4) at (4,25) {BB4\\\texttt{1000}};
  \node[bb, fill=ln-source!40] (b5) at (20,25) {BB5\\\texttt{0100}};
  \node[bb, fill=ln-accent!20] (b6) at (22,9) {BB6\\\texttt{1111}};
  \draw[->] (b1) -- (b2);
  \draw[->] (b1) -- (b3);
  \draw[->] (b2) -- (b4);
  \draw[->] (b2) -- (b5);
  \draw[->] (b4) -- (b6);
  \draw[->] (b5) -- (b6);
  \draw[->] (b3) |- (b6);
  % stack snapshots: \snap{x}{y}{title}{list of alt/reconv}
  \newcommand{\snapshot}[4]{%
    \node[anchor=south west, font=\scriptsize\sffamily\bfseries] at (#1,#2+1) {#3};
    \node[anchor=south, text=ln-muted] at (#1+8,#2-4.5) {\iflangja{代替 PC}{alt.\ PC}};
    \node[anchor=south, text=ln-muted] at (#1+21.5,#2-4.5) {\iflangja{再収束点}{reconv.}};
    \foreach \alt/\rec [count=\k from 0] in {#4} {
      \draw[fill=ln-box] (#1,#2-5-\k*6) rectangle ++(16,-5.5)
        node[midway] {\alt};
      \draw[fill=ln-box] (#1+16,#2-5-\k*6) rectangle ++(11,-5.5)
        node[midway] {\rec};
    }}
  \snapshot{52}{56}{\iflangja{(a) BB1 の後}{(a) after BB1}}
    {{BB3 \texttt{0011}}/BB6}
  \snapshot{84}{56}{\iflangja{(b) BB2 の後}{(b) after BB2}}
    {{BB5 \texttt{0100}}/BB6, {BB3 \texttt{0011}}/BB6}
  \snapshot{52}{26}{\iflangja{(c) BB4 の後（pop）}{(c) after BB4 (pop)}}
    {{BB3 \texttt{0011}}/BB6}
  \node[anchor=south west, font=\scriptsize\sffamily\bfseries] at (84,27)
    {\iflangja{(d) BB5 の後（pop）}{(d) after BB5 (pop)}};
  \node[anchor=north west, text=ln-muted, align=left] at (84,24)
    {\iflangja{空：BB3 の後は\\BB6 を \texttt{1111} で実行}{empty: after BB3,\\BB6 runs with \texttt{1111}}};
\end{tikzpicture}

% Width: body (120 mm); the drawing is about 112 mm wide.
\begin{tikzpicture}[x=1mm, y=1mm, font=\scriptsize\sffamily,
    >={Stealth[length=1.6mm]},
    bb/.style={draw, minimum width=13mm, minimum height=7.5mm,
               align=center, inner sep=1pt}]
  % nested CFG
  \node[bb, fill=ln-accent!20] (b1) at (22,56) {BB1\\\texttt{1111}};
  \node[bb, fill=ln-source!25] (b2) at (12,41) {BB2\\\texttt{1100}};
  \node[bb, fill=ln-tip!25] (b3) at (36,41) {BB3\\\texttt{0011}};
  \node[bb, fill=ln-source!12] (b4) at (4,25) {BB4\\\texttt{1000}};
  \node[bb, fill=ln-source!40] (b5) at (20,25) {BB5\\\texttt{0100}};
  \node[bb, fill=ln-accent!20] (b6) at (22,9) {BB6\\\texttt{1111}};
  \draw[->] (b1) -- (b2);
  \draw[->] (b1) -- (b3);
  \draw[->] (b2) -- (b4);
  \draw[->] (b2) -- (b5);
  \draw[->] (b4) -- (b6);
  \draw[->] (b5) -- (b6);
  \draw[->] (b3) |- (b6);
  % stack snapshots: \snap{x}{y}{title}{list of alt/reconv}
  \newcommand{\snapshot}[4]{%
    \node[anchor=south west, font=\scriptsize\sffamily\bfseries] at (#1,#2+1) {#3};
    \node[anchor=south, text=ln-muted] at (#1+8,#2-4.5) {\iflangja{代替 PC}{alt.\ PC}};
    \node[anchor=south, text=ln-muted] at (#1+21.5,#2-4.5) {\iflangja{再収束点}{reconv.}};
    \foreach \alt/\rec [count=\k from 0] in {#4} {
      \draw[fill=ln-box] (#1,#2-5-\k*6) rectangle ++(16,-5.5)
        node[midway] {\alt};
      \draw[fill=ln-box] (#1+16,#2-5-\k*6) rectangle ++(11,-5.5)
        node[midway] {\rec};
    }}
  \snapshot{52}{56}{\iflangja{(a) BB1 の後}{(a) after BB1}}
    {{BB3 \texttt{0011}}/BB6}
  \snapshot{84}{56}{\iflangja{(b) BB2 の後}{(b) after BB2}}
    {{BB5 \texttt{0100}}/BB6, {BB3 \texttt{0011}}/BB6}
  \snapshot{52}{26}{\iflangja{(c) BB4 の後（pop）}{(c) after BB4 (pop)}}
    {{BB3 \texttt{0011}}/BB6}
  \node[anchor=south west, font=\scriptsize\sffamily\bfseries] at (84,27)
    {\iflangja{(d) BB5 の後（pop）}{(d) after BB5 (pop)}};
  \node[anchor=north west, text=ln-muted, align=left] at (84,24)
    {\iflangja{空：BB3 の後は\\BB6 を \texttt{1111} で実行}{empty: after BB3,\\BB6 runs with \texttt{1111}}};
\end{tikzpicture}

% Width: body (120 mm); the drawing is about 112 mm wide.
\begin{tikzpicture}[x=1mm, y=1mm, font=\scriptsize\sffamily,
    >={Stealth[length=1.6mm]},
    bb/.style={draw, minimum width=13mm, minimum height=7.5mm,
               align=center, inner sep=1pt}]
  % nested CFG
  \node[bb, fill=ln-accent!20] (b1) at (22,56) {BB1\\\texttt{1111}};
  \node[bb, fill=ln-source!25] (b2) at (12,41) {BB2\\\texttt{1100}};
  \node[bb, fill=ln-tip!25] (b3) at (36,41) {BB3\\\texttt{0011}};
  \node[bb, fill=ln-source!12] (b4) at (4,25) {BB4\\\texttt{1000}};
  \node[bb, fill=ln-source!40] (b5) at (20,25) {BB5\\\texttt{0100}};
  \node[bb, fill=ln-accent!20] (b6) at (22,9) {BB6\\\texttt{1111}};
  \draw[->] (b1) -- (b2);
  \draw[->] (b1) -- (b3);
  \draw[->] (b2) -- (b4);
  \draw[->] (b2) -- (b5);
  \draw[->] (b4) -- (b6);
  \draw[->] (b5) -- (b6);
  \draw[->] (b3) |- (b6);
  % stack snapshots: \snap{x}{y}{title}{list of alt/reconv}
  \newcommand{\snapshot}[4]{%
    \node[anchor=south west, font=\scriptsize\sffamily\bfseries] at (#1,#2+1) {#3};
    \node[anchor=south, text=ln-muted] at (#1+8,#2-4.5) {\iflangja{代替 PC}{alt.\ PC}};
    \node[anchor=south, text=ln-muted] at (#1+21.5,#2-4.5) {\iflangja{再収束点}{reconv.}};
    \foreach \alt/\rec [count=\k from 0] in {#4} {
      \draw[fill=ln-box] (#1,#2-5-\k*6) rectangle ++(16,-5.5)
        node[midway] {\alt};
      \draw[fill=ln-box] (#1+16,#2-5-\k*6) rectangle ++(11,-5.5)
        node[midway] {\rec};
    }}
  \snapshot{52}{56}{\iflangja{(a) BB1 の後}{(a) after BB1}}
    {{BB3 \texttt{0011}}/BB6}
  \snapshot{84}{56}{\iflangja{(b) BB2 の後}{(b) after BB2}}
    {{BB5 \texttt{0100}}/BB6, {BB3 \texttt{0011}}/BB6}
  \snapshot{52}{26}{\iflangja{(c) BB4 の後（pop）}{(c) after BB4 (pop)}}
    {{BB3 \texttt{0011}}/BB6}
  \node[anchor=south west, font=\scriptsize\sffamily\bfseries] at (84,27)
    {\iflangja{(d) BB5 の後（pop）}{(d) after BB5 (pop)}};
  \node[anchor=north west, text=ln-muted, align=left] at (84,24)
    {\iflangja{空：BB3 の後は\\BB6 を \texttt{1111} で実行}{empty: after BB3,\\BB6 runs with \texttt{1111}}};
\end{tikzpicture}

  \caption{The SIMT stack on a nested branch.  Left: BB2 diverges again
    and all paths reconverge at BB6.  Right: the stack (top entry
    first) after the branch in BB1, after the branch in BB2, after the
    pop at the end of BB4 and after the pop at the end of BB5.}
  \label{fig:l06-simt-stack}
\end{figure}

\Cref{tab:l06-simt-trace} traces the nested example.  BB4 ends with the
next PC BB6, which equals the reconvergence PC on top of the stack, so
the warp pops and fetches BB5 with \texttt{0100}; BB5 ends at BB6 again,
the warp pops the second entry and fetches BB3 with \texttt{0011}; when
BB3 ends the stack is empty and BB6 runs with the full mask.  The
two-field stack does not say where \texttt{1111} comes from; the
three-field stack of the sidenote restores it from the pre-branch
entry.  Every
thread has executed exactly its own path, and the six blocks took six
fetches.\sidenote{In the three-field formulation of K12 §1.4.2 and A18
ch.~3 \cite{kim2012gpgpu,aamodt2018gpu} the top of the stack always
describes what the warp is executing.  At a divergent branch the top
entry's next PC is replaced by the reconvergence PC (keeping the
pre-branch mask), and one entry per path is pushed, each with its own
next PC, mask and the common reconvergence PC.  When the warp reaches
the reconvergence PC of the top entry, that entry is popped and the
next one supplies PC and mask, so the pre-branch entry restores
\texttt{1111} at BB6.}

\begin{table}[tb]
  \centering
  \caption{Execution of the nested branch of \cref{fig:l06-simt-stack}
    with a SIMT stack.  Stack entries are written
    (alternate PC, mask; reconvergence PC), top first.}
  \label{tab:l06-simt-trace}
  \setlength{\tabcolsep}{4pt}
  \begin{tabular}{@{}l l l >{\raggedright\arraybackslash}p{30mm} >{\raggedright\arraybackslash}p{34mm}@{}}
    \toprule
    block & mask & next PC & action & stack afterwards \\
    \midrule
    BB1 & \texttt{1111} & BB2 / BB3 & diverges: push BB3, go to BB2
      & (BB3, \texttt{0011}; BB6) \\
    BB2 & \texttt{1100} & BB4 / BB5 & diverges: push BB5, go to BB4
      & (BB5, \texttt{0100}; BB6), (BB3, \texttt{0011}; BB6) \\
    BB4 & \texttt{1000} & BB6 & equals top reconv.: pop, fetch BB5
      & (BB3, \texttt{0011}; BB6) \\
    BB5 & \texttt{0100} & BB6 & equals top reconv.: pop, fetch BB3
      & empty \\
    BB3 & \texttt{0011} & BB6 & stack empty: full mask (implicit) & empty \\
    BB6 & \texttt{1111} & --- & --- & empty \\
    \bottomrule
  \end{tabular}
\end{table}

The stack has a hardware cost of a few entries per warp and a
reconvergence comparison per fetch; the number of entries it must
hold grows with the nesting depth of divergent branches.

\subsection{Finding the reconvergence point}

The reconvergence PC is computed at compile time: the compiler builds
the control-flow graph, finds the immediate post-dominator of every
branch, and either marks it in the code or inserts explicit
instructions that push and pop the stack.  The stack itself is kept in hardware, because only
at run time is it known which branches actually diverge.  The compiler
could in principle avoid the stack operations for branches that never
diverge, but deciding statically whether a branch diverges is hard, so
a compiler must treat a branch as divergent whenever it cannot prove
otherwise, and over-reports divergent branches.\margin{Lesson~11
presents divergence analysis, the compiler analysis that decides which
values, and hence which branch conditions, can differ between the
threads of a warp.}

\begin{note}
  The SIMT stack serialises the paths of a branch: the warp runs one
  path to the reconvergence point before it starts the other.  A thread
  that spins on a value written by a thread on the other path hangs if
  its own path is scheduled first, since the writer runs only after the
  spinning path reaches the reconvergence point, which it never does:
  a \emph{SIMT deadlock}.  For example, lane~0 takes the \texttt{then}
  path and loops until a flag in memory is set, while lane~1 takes the
  \texttt{else} path, which sets the flag.  If the stack runs the
  \texttt{then} path first, lane~0 loops forever and the \texttt{else}
  path is never fetched; a spin lock taken by several lanes of one warp
  fails the same way.  A18 ch.~3 discusses the problem together with
  stackless designs that keep a program counter per thread, so that the
  scheduler can switch to the other path.\margin{Since Volta, NVIDIA
  GPUs keep a program counter and call stack per thread (independent
  thread scheduling, Lesson~5) and no longer depend on the classic
  stack for correctness (CPG §4.1 SIMT Architecture
  \cite{nvidia2024cudaguide}).}  The rest of this book keeps the
  per-warp stack model of this section.
\end{note}

\subsection{Large warps}

The warp size is a microarchitectural choice, not a law: one fetch and
decode can drive 32, 64 or 128 lanes, and a very wide warp turns the SM
into a wide vector machine with little front-end cost per
lane.  The cost
is divergence: the more threads a warp holds, the more likely it is
that one of them disagrees with the others, and every divergence
serialises paths and occupies the SIMT stack.  The right width depends
on how uniform the application's control flow is.\caution{Real branch outcomes are correlated (conditions on
\texttt{threadIdx.x / 32} are uniform per warp by construction), so
independent outcomes are a pessimistic model; it shows the trend, not
a real divergence rate.}

\begin{example}
  \label{ex:l06-large-warp}
  Suppose each thread takes a branch independently with probability
  $p = 0.1$.  A warp of $W$ threads is uniform only if all threads
  agree, which happens with probability $0.9^W + 0.1^W$: about
  $0.66$ for $W = 4$, $0.034$ for $W = 32$ and $0.0012$ for $W = 64$.
  Under this model almost every branch of a 32-thread warp
  diverges.
\end{example}

\begin{definition}[Large warp]
  \label{def:l06-large-warp}
  A \term{large warp}[ラージワープ] is a warp with more threads than
  the conventional 32 (64, 128 or more), whose instructions are
  fetched and decoded once for all of them, so that the front-end cost
  is shared by more lanes.
\end{definition}

\caution{Lecture vs paper: this chapter, following the lecture,
treats a large warp as a wider warp.  Narasiman et al.\ \cite{narasiman2011twolevel} pack the
\emph{active} threads of a large warp into SIMD-wide sub-warps every
cycle, so divergence leaves fewer idle lanes; combined with two-level
warp scheduling (Lesson~7) they report 19.1\,\% over traditional GPU
cores.}%
The large warp of the paper is thus also a way of fighting divergence,
which leads to the next idea.

\subsection{Dynamic warp formation}

After a divergent branch each path runs with only part of its lanes
busy; in \cref{fig:l06-divergent-cfg} half of the lanes idle in BB2 and
in BB3.  The
fraction of busy lanes is the measure of this loss.

\begin{definition}[SIMD efficiency]
  \label{def:l06-simd-efficiency}
  The \term{SIMD efficiency}[SIMD効率] of an execution is the number of
  lane-slots that do work divided by the number of lane-slots issued
  (warp size $\times$ warp instructions issued).  A program without
  divergence has a SIMD efficiency of 100\,\%; the divergent warp of
  \cref{ex:l06-if-conversion} has $40/64 = 62.5\,\%$.
\end{definition}
\margin{PMPP ch.~10 \cite{hwu2022pmpp} applies the same measure in
software: it rewrites a parallel reduction so that the active threads
stay packed in few warps, keeping SIMD efficiency high.}

The many warps of an SM are often at the same point of the program, and
after the same divergent branch, different warps are waiting at the
same PC with complementary masks.  Their threads can be combined.\margin{Fung et al.\ report an average speed-up
of 20.7\,\% for an estimated area increase of 4.7\,\% from a
realistic DWF implementation \cite{fung2007dwf}.}

\begin{definition}[Dynamic warp formation]
  \label{def:l06-dwf}
  The technique of \term{dynamic warp formation (DWF)}[動的ワープ形成]\index{DWF|see{dynamic warp formation (DWF)}}
  regroups, at run
  time, threads of different warps that are about to execute the same PC
  into new warps with fuller active masks
  (\cref{fig:l06-dwf}).
\end{definition}

\begin{figure}[tb]
  \centering
  % Dynamic warp formation (Lesson 6).  Two warps of four threads (T0-T3,
% T4-T7) diverge on the same branch: at BB2 warp 0 has lanes 0-1 active
% and warp 1 lanes 2-3, so the two half warps fill one new warp without
% moving any thread out of its home lane; BB3 likewise.
% Redraws the outline figure m6l1_dynamicwarp.
% Usage: % Dynamic warp formation (Lesson 6).  Two warps of four threads (T0-T3,
% T4-T7) diverge on the same branch: at BB2 warp 0 has lanes 0-1 active
% and warp 1 lanes 2-3, so the two half warps fill one new warp without
% moving any thread out of its home lane; BB3 likewise.
% Redraws the outline figure m6l1_dynamicwarp.
% Usage: % Dynamic warp formation (Lesson 6).  Two warps of four threads (T0-T3,
% T4-T7) diverge on the same branch: at BB2 warp 0 has lanes 0-1 active
% and warp 1 lanes 2-3, so the two half warps fill one new warp without
% moving any thread out of its home lane; BB3 likewise.
% Redraws the outline figure m6l1_dynamicwarp.
% Usage: \input{figures/l06-dwf}
% Width: body (120 mm); the drawing is about 114 mm wide.
\begin{tikzpicture}[x=1mm, y=1mm, font=\scriptsize\sffamily,
    >={Stealth[length=1.6mm]}]
  % \lanes{x}{y}{labels}{fill}: four lane cells, "-" = idle
  \def\dashmark{-}
  \newcommand{\lanes}[4]{%
    \foreach \t [count=\l from 0] in {#3} {
      \ifx\t\dashmark
        \draw[ln-rule, dashed] (#1+\l*7,#2) rectangle ++(6,5);
      \else
        \draw[fill=#4] (#1+\l*7,#2) rectangle ++(6,5) node[midway] {\t};
      \fi
    }}
  % column heads
  \node[font=\scriptsize\sffamily\bfseries] at (13.5,30) {\iflangja{ワープ 0}{warp 0}};
  \node[font=\scriptsize\sffamily\bfseries] at (49.5,30) {\iflangja{ワープ 1}{warp 1}};
  \node[font=\scriptsize\sffamily\bfseries, align=center] at (95.5,30)
    {\iflangja{動的に形成したワープ}{dynamically formed warp}};
  \foreach \l in {0,...,3} {
    \node[text=ln-muted] at (3+\l*7,25.5) {\l};
    \node[text=ln-muted] at (39+\l*7,25.5) {\l};
    \node[text=ln-muted] at (85+\l*7,25.5) {\l};
  }
  \node[anchor=east, text=ln-muted] at (-1,25.5) {\iflangja{レーン}{lane}};
  % BB2 row
  \node[anchor=east] at (-1,19.5) {BB2};
  \lanes{0}{17}{T0,T1,-,-}{ln-source!25}
  \lanes{36}{17}{-,-,T6,T7}{ln-source!25}
  \lanes{82}{17}{T0,T1,T6,T7}{ln-source!25}
  % BB3 row
  \node[anchor=east] at (-1,9.5) {BB3};
  \lanes{0}{7}{-,-,T2,T3}{ln-tip!25}
  \lanes{36}{7}{T4,T5,-,-}{ln-tip!25}
  \lanes{82}{7}{T4,T5,T2,T3}{ln-tip!25}
  \draw[->, very thick, ln-accent] (66,14.5) -- (79,14.5);
  \node[anchor=north, text=ln-muted, align=center] at (72.5,12.5)
    {\iflangja{同じ PC で\\再グループ化}{regroup\\at equal PC}};
  \node[anchor=north west, text=ln-muted] at (0,4)
    {\iflangja{SIMD 効率 50\,\% $\times$ 4 命令スロット}{SIMD efficiency 50\,\% in 4 issue slots}};
  \node[anchor=north west, text=ln-muted] at (82,4)
    {\iflangja{100\,\% $\times$ 2 スロット}{100\,\% in 2 issue slots}};
\end{tikzpicture}

% Width: body (120 mm); the drawing is about 114 mm wide.
\begin{tikzpicture}[x=1mm, y=1mm, font=\scriptsize\sffamily,
    >={Stealth[length=1.6mm]}]
  % \lanes{x}{y}{labels}{fill}: four lane cells, "-" = idle
  \def\dashmark{-}
  \newcommand{\lanes}[4]{%
    \foreach \t [count=\l from 0] in {#3} {
      \ifx\t\dashmark
        \draw[ln-rule, dashed] (#1+\l*7,#2) rectangle ++(6,5);
      \else
        \draw[fill=#4] (#1+\l*7,#2) rectangle ++(6,5) node[midway] {\t};
      \fi
    }}
  % column heads
  \node[font=\scriptsize\sffamily\bfseries] at (13.5,30) {\iflangja{ワープ 0}{warp 0}};
  \node[font=\scriptsize\sffamily\bfseries] at (49.5,30) {\iflangja{ワープ 1}{warp 1}};
  \node[font=\scriptsize\sffamily\bfseries, align=center] at (95.5,30)
    {\iflangja{動的に形成したワープ}{dynamically formed warp}};
  \foreach \l in {0,...,3} {
    \node[text=ln-muted] at (3+\l*7,25.5) {\l};
    \node[text=ln-muted] at (39+\l*7,25.5) {\l};
    \node[text=ln-muted] at (85+\l*7,25.5) {\l};
  }
  \node[anchor=east, text=ln-muted] at (-1,25.5) {\iflangja{レーン}{lane}};
  % BB2 row
  \node[anchor=east] at (-1,19.5) {BB2};
  \lanes{0}{17}{T0,T1,-,-}{ln-source!25}
  \lanes{36}{17}{-,-,T6,T7}{ln-source!25}
  \lanes{82}{17}{T0,T1,T6,T7}{ln-source!25}
  % BB3 row
  \node[anchor=east] at (-1,9.5) {BB3};
  \lanes{0}{7}{-,-,T2,T3}{ln-tip!25}
  \lanes{36}{7}{T4,T5,-,-}{ln-tip!25}
  \lanes{82}{7}{T4,T5,T2,T3}{ln-tip!25}
  \draw[->, very thick, ln-accent] (66,14.5) -- (79,14.5);
  \node[anchor=north, text=ln-muted, align=center] at (72.5,12.5)
    {\iflangja{同じ PC で\\再グループ化}{regroup\\at equal PC}};
  \node[anchor=north west, text=ln-muted] at (0,4)
    {\iflangja{SIMD 効率 50\,\% $\times$ 4 命令スロット}{SIMD efficiency 50\,\% in 4 issue slots}};
  \node[anchor=north west, text=ln-muted] at (82,4)
    {\iflangja{100\,\% $\times$ 2 スロット}{100\,\% in 2 issue slots}};
\end{tikzpicture}

% Width: body (120 mm); the drawing is about 114 mm wide.
\begin{tikzpicture}[x=1mm, y=1mm, font=\scriptsize\sffamily,
    >={Stealth[length=1.6mm]}]
  % \lanes{x}{y}{labels}{fill}: four lane cells, "-" = idle
  \def\dashmark{-}
  \newcommand{\lanes}[4]{%
    \foreach \t [count=\l from 0] in {#3} {
      \ifx\t\dashmark
        \draw[ln-rule, dashed] (#1+\l*7,#2) rectangle ++(6,5);
      \else
        \draw[fill=#4] (#1+\l*7,#2) rectangle ++(6,5) node[midway] {\t};
      \fi
    }}
  % column heads
  \node[font=\scriptsize\sffamily\bfseries] at (13.5,30) {\iflangja{ワープ 0}{warp 0}};
  \node[font=\scriptsize\sffamily\bfseries] at (49.5,30) {\iflangja{ワープ 1}{warp 1}};
  \node[font=\scriptsize\sffamily\bfseries, align=center] at (95.5,30)
    {\iflangja{動的に形成したワープ}{dynamically formed warp}};
  \foreach \l in {0,...,3} {
    \node[text=ln-muted] at (3+\l*7,25.5) {\l};
    \node[text=ln-muted] at (39+\l*7,25.5) {\l};
    \node[text=ln-muted] at (85+\l*7,25.5) {\l};
  }
  \node[anchor=east, text=ln-muted] at (-1,25.5) {\iflangja{レーン}{lane}};
  % BB2 row
  \node[anchor=east] at (-1,19.5) {BB2};
  \lanes{0}{17}{T0,T1,-,-}{ln-source!25}
  \lanes{36}{17}{-,-,T6,T7}{ln-source!25}
  \lanes{82}{17}{T0,T1,T6,T7}{ln-source!25}
  % BB3 row
  \node[anchor=east] at (-1,9.5) {BB3};
  \lanes{0}{7}{-,-,T2,T3}{ln-tip!25}
  \lanes{36}{7}{T4,T5,-,-}{ln-tip!25}
  \lanes{82}{7}{T4,T5,T2,T3}{ln-tip!25}
  \draw[->, very thick, ln-accent] (66,14.5) -- (79,14.5);
  \node[anchor=north, text=ln-muted, align=center] at (72.5,12.5)
    {\iflangja{同じ PC で\\再グループ化}{regroup\\at equal PC}};
  \node[anchor=north west, text=ln-muted] at (0,4)
    {\iflangja{SIMD 効率 50\,\% $\times$ 4 命令スロット}{SIMD efficiency 50\,\% in 4 issue slots}};
  \node[anchor=north west, text=ln-muted] at (82,4)
    {\iflangja{100\,\% $\times$ 2 スロット}{100\,\% in 2 issue slots}};
\end{tikzpicture}

  \caption{Dynamic warp formation.  Warps 0 and 1 diverge on the same
    branch; at BB2 the active threads of warp~0 sit in lanes 0--1 and
    those of warp~1 in lanes 2--3, so they form one full warp without
    leaving their lanes.  BB3 is merged the same way.}
  \label{fig:l06-dwf}
\end{figure}

The constraint is the register file.  Lesson~4 assigned the registers
of thread~$k$ to bank~$k$, so that lane~$k$ reads from its own bank; a
thread moved into another lane
would need its registers from a bank that belongs to another lane,
creating the register bank conflicts the layout was designed to
avoid.  DWF therefore
keeps every thread in its home lane and merges only warps whose active
lanes do not overlap, as in \cref{fig:l06-dwf};\tip{DWF can merge two warps at the
same PC only if the bitwise AND of their masks is zero; the new mask is
their OR.} two half warps that both
use lanes 0--1 cannot be merged.

\begin{example}
  \label{ex:l06-dwf}
  Take two warps executing a branch with the block sizes of
  \cref{ex:l06-if-conversion}, whose condition depends on the data so
  that warp~0 takes BB2 in lanes 0--1 and warp~1 in lanes 2--3, as in
  \cref{fig:l06-dwf}.  Without DWF, BB2 and BB3 cost
  $2 \times (6+6) = 24$ issue slots at 50\,\% SIMD efficiency.  With DWF
  the merged warps need $6+6 = 12$ issue slots at 100\,\%.  The whole
  kernel drops from $2 \times 16 = 32$ to $2\cdot2 + 12 + 2\cdot2 = 20$
  warp instructions, and its SIMD efficiency rises from 62.5\,\% to
  $80/80 = 100\,\%$.
\end{example}

\begin{keypoint}
  A divergent branch makes a warp run both paths under partial masks.
  If-conversion trades it for always fetching both paths; the SIMT
  stack runs the paths one after the other and reconverges at the
  immediate post-dominator.  Wider warps share the front end but
  diverge more (the large warp microarchitecture of Narasiman et al.\
  repacks active threads into sub-warps), and dynamic warp formation
  recovers idle lanes, keeping each thread in its home lane.
\end{keypoint}

\section{Register optimizations}
\label{sec:l06-registers}

The register file of an SM holds a private register set for every
resident thread (Lesson~4), so its size is threads times registers per
thread: \SI{64}{\kibi\byte}, \SI{128}{\kibi\byte} or
\SI{256}{\kibi\byte} per SM in successive
generations.\source{M6, register optimizations; K12 §1.4.1 GPU
Pipeline; \cite{gebhart2011energy,gebhart2011compile,jeon2015rfv}.}
A large SRAM is slow and consumes power, and a GPU cannot accept a slow
one: the register file must deliver operands at the rate at which the
lanes consume them, so its read bandwidth must equal the execution
bandwidth.\margin{\SI{256}{\kibi\byte} per SM on V100, A100 and H100;
per GPU \SI{27648}{\kibi\byte} on A100 and \SI{33792}{\kibi\byte} on
H100 SXM5 (preliminary figure) \cite[Table~3]{nvidia2022hopper}.}  Two remedies come from CPU design: a
hierarchy, to make the common access fast, and sharing, to need less
storage.  Four observations about how GPU programs use registers show
where they apply.

\subsection{Four observations}

\begin{description}
  \item[Not all threads are active.]  Threads of the same kernel use the
    same register names, but they do not finish at the same time:
    execution is asynchronous, and a thread that has finished still
    occupies its registers until its warp or block is released.  On a
    divergent branch, the two paths may also use different registers.
  \item[Not all registers are live.]  In
    \texttt{add r1, r2, \#1} followed by
    \texttt{mov r2, \#10}, the value in \texttt{r2} is
    never read again after the \texttt{add}: it is dead until the
    \texttt{mov} overwrites it.
  \item[Producer and consumer are close.]  In
    \texttt{mov r2, \#10} followed by
    \texttt{add r3, r2, \#20}, the value is consumed by
    the very next instruction, and a scheduler can keep the two
    together.  Such a value may not need to reach the large register
    file at all.
  \item[Most values are read only once.]  The first value of
    \texttt{r2} in \texttt{mov r2, \#10} followed by
    \texttt{add r2, r2, \#20} is read by one instruction and then
    overwritten.  Many
    values follow this pattern; storing them in the main register file
    is questionable.\caution{Lecture vs paper: the lecture gives almost
    70\,\% of register values as read only once.  Gebhart et al.\
    measure \enquote{up to 70\,\%} across their workloads, with HPC
    codes reusing more (40\,\% of values read more than once), and
    \enquote{up to 40\,\%} read once \emph{and} within 3 instructions
    of being produced \cite{gebhart2011energy}.}
\end{description}

\begin{definition}[Live and dead registers]
  \label{def:l06-live-dead}
  A register is a \term{live register}[生存レジスタ] at a point of the
  program if the value it holds will be read later before being
  overwritten, and a \term{dead register}[非生存レジスタ] otherwise.
  The \term{value lifetime}[値の生存期間] of a value is the interval
  from the instruction that produces it to its last read.
\end{definition}

The second observation weighs more on a GPU than on a CPU
(\cref{fig:l06-live-dead}).  On a CPU, the \texttt{add} and the
\texttt{mov} execute in consecutive cycles, and \texttt{r2} is dead for
a cycle or less.  On a GPU the scheduler interleaves warps, so the next
instruction of the same warp may issue tens of warp instructions later;
for that whole time the dead value of \texttt{r2} occupies a register
of every thread of the warp.\margin{Lesson~7 discusses the warp
schedulers that decide this interleaving; a greedy scheduler shortens
the gap, a round-robin one stretches it to the number of ready warps.}

\begin{figure}[tb]
  \centering
  % Live and dead registers on a CPU and on a GPU (Lesson 6).  After
% "add r1, r2, #1" register r2 is dead until "mov r2, #10" overwrites it.
% On a CPU the two instructions issue in consecutive cycles; on a GPU the
% issue slots in between go to other warps, so r2 of warp w1 holds a
% dead value for tens of cycles.
% Redraws the outline figure m6l2_o2.
% Usage: % Live and dead registers on a CPU and on a GPU (Lesson 6).  After
% "add r1, r2, #1" register r2 is dead until "mov r2, #10" overwrites it.
% On a CPU the two instructions issue in consecutive cycles; on a GPU the
% issue slots in between go to other warps, so r2 of warp w1 holds a
% dead value for tens of cycles.
% Redraws the outline figure m6l2_o2.
% Usage: % Live and dead registers on a CPU and on a GPU (Lesson 6).  After
% "add r1, r2, #1" register r2 is dead until "mov r2, #10" overwrites it.
% On a CPU the two instructions issue in consecutive cycles; on a GPU the
% issue slots in between go to other warps, so r2 of warp w1 holds a
% dead value for tens of cycles.
% Redraws the outline figure m6l2_o2.
% Usage: \input{figures/l06-live-dead}
% Width: body (120 mm); the drawing is about 116 mm wide.
\begin{tikzpicture}[x=1mm, y=1mm, font=\scriptsize\sffamily,
    slot/.style={draw, minimum width=14mm, minimum height=6mm,
                 inner sep=1pt, align=center}]
  % CPU
  \node[anchor=west, font=\scriptsize\sffamily\bfseries] at (0,31) {(a) CPU};
  \node[slot, fill=ln-accent!20] at (24,31) {\texttt{add r1,r2,\#1}};
  \node[slot, fill=ln-tip!25] at (40,31) {\texttt{mov r2,\#10}};
  \draw[ln-caution, thick] (31,26.5) -- ++(0,-1) -- ++(2,0) -- ++(0,1);
  \node[anchor=west, text=ln-caution] at (34,24.5)
    {\iflangja{\texttt{r2} 非生存：1 サイクル未満}{\texttt{r2} dead: under one cycle}};
  % GPU
  \node[anchor=west, font=\scriptsize\sffamily\bfseries] at (0,13) {(b) GPU};
  \foreach \t/\w/\c [count=\i from 0] in
      {add/w1/ln-accent!20, add/w2/ln-box, add/w3/ln-box, add/w4/ln-box,
       {\ldots}/{}/white, mov/w1/ln-tip!25} {
    \node[slot, fill=\c, minimum width=13mm] at (24+\i*15,13) {\texttt{\t}\\\w};
  }
  \draw[ln-caution, thick] (31,8.5) -- ++(0,-2) -- ++(72,0) -- ++(0,2);
  \node[anchor=north, text=ln-caution] at (67,6)
    {\iflangja{w1 の \texttt{r2} は非生存（数十命令スロットの間）}{\texttt{r2} of w1 dead for tens of issue slots}};
  \node[anchor=north, text=ln-muted] at (67,20.5)
    {\iflangja{1 命令スロットごとに 1 ワープ命令}{one warp instruction per issue slot}};
\end{tikzpicture}

% Width: body (120 mm); the drawing is about 116 mm wide.
\begin{tikzpicture}[x=1mm, y=1mm, font=\scriptsize\sffamily,
    slot/.style={draw, minimum width=14mm, minimum height=6mm,
                 inner sep=1pt, align=center}]
  % CPU
  \node[anchor=west, font=\scriptsize\sffamily\bfseries] at (0,31) {(a) CPU};
  \node[slot, fill=ln-accent!20] at (24,31) {\texttt{add r1,r2,\#1}};
  \node[slot, fill=ln-tip!25] at (40,31) {\texttt{mov r2,\#10}};
  \draw[ln-caution, thick] (31,26.5) -- ++(0,-1) -- ++(2,0) -- ++(0,1);
  \node[anchor=west, text=ln-caution] at (34,24.5)
    {\iflangja{\texttt{r2} 非生存：1 サイクル未満}{\texttt{r2} dead: under one cycle}};
  % GPU
  \node[anchor=west, font=\scriptsize\sffamily\bfseries] at (0,13) {(b) GPU};
  \foreach \t/\w/\c [count=\i from 0] in
      {add/w1/ln-accent!20, add/w2/ln-box, add/w3/ln-box, add/w4/ln-box,
       {\ldots}/{}/white, mov/w1/ln-tip!25} {
    \node[slot, fill=\c, minimum width=13mm] at (24+\i*15,13) {\texttt{\t}\\\w};
  }
  \draw[ln-caution, thick] (31,8.5) -- ++(0,-2) -- ++(72,0) -- ++(0,2);
  \node[anchor=north, text=ln-caution] at (67,6)
    {\iflangja{w1 の \texttt{r2} は非生存（数十命令スロットの間）}{\texttt{r2} of w1 dead for tens of issue slots}};
  \node[anchor=north, text=ln-muted] at (67,20.5)
    {\iflangja{1 命令スロットごとに 1 ワープ命令}{one warp instruction per issue slot}};
\end{tikzpicture}

% Width: body (120 mm); the drawing is about 116 mm wide.
\begin{tikzpicture}[x=1mm, y=1mm, font=\scriptsize\sffamily,
    slot/.style={draw, minimum width=14mm, minimum height=6mm,
                 inner sep=1pt, align=center}]
  % CPU
  \node[anchor=west, font=\scriptsize\sffamily\bfseries] at (0,31) {(a) CPU};
  \node[slot, fill=ln-accent!20] at (24,31) {\texttt{add r1,r2,\#1}};
  \node[slot, fill=ln-tip!25] at (40,31) {\texttt{mov r2,\#10}};
  \draw[ln-caution, thick] (31,26.5) -- ++(0,-1) -- ++(2,0) -- ++(0,1);
  \node[anchor=west, text=ln-caution] at (34,24.5)
    {\iflangja{\texttt{r2} 非生存：1 サイクル未満}{\texttt{r2} dead: under one cycle}};
  % GPU
  \node[anchor=west, font=\scriptsize\sffamily\bfseries] at (0,13) {(b) GPU};
  \foreach \t/\w/\c [count=\i from 0] in
      {add/w1/ln-accent!20, add/w2/ln-box, add/w3/ln-box, add/w4/ln-box,
       {\ldots}/{}/white, mov/w1/ln-tip!25} {
    \node[slot, fill=\c, minimum width=13mm] at (24+\i*15,13) {\texttt{\t}\\\w};
  }
  \draw[ln-caution, thick] (31,8.5) -- ++(0,-2) -- ++(72,0) -- ++(0,2);
  \node[anchor=north, text=ln-caution] at (67,6)
    {\iflangja{w1 の \texttt{r2} は非生存（数十命令スロットの間）}{\texttt{r2} of w1 dead for tens of issue slots}};
  \node[anchor=north, text=ln-muted] at (67,20.5)
    {\iflangja{1 命令スロットごとに 1 ワープ命令}{one warp instruction per issue slot}};
\end{tikzpicture}

  \caption{After \texttt{add r1,r2,\#1}, register \texttt{r2} is dead
    until \texttt{mov r2,\#10}.  (a) On a CPU the two instructions are
    adjacent.  (b) On a GPU the issue slots in between belong to other
    warps, so the dead interval of warp \texttt{w1}'s \texttt{r2}
    spans tens of instructions.}
  \label{fig:l06-live-dead}
\end{figure}


\begin{example}
  \label{ex:l06-live-dead}
  An SM interleaves 24 warps round-robin, one instruction each per round,
  so consecutive instructions of one warp are 24 cycles apart.  A value
  produced by instruction $k$, last read by instruction $k+1$ and
  overwritten by instruction $k+5$ is live for $24$ cycles and dead for
  $4 \times 24 = 96$: its register holds useful data
  $24/120 = 20\,\%$ of the time.  On a CPU issuing one instruction per
  cycle the same value is live for 1 cycle and dead for 4, with the same
  20\,\%, but the dead register costs one entry of a small file; on the
  GPU it costs one entry per thread for 96 cycles in a file of hundreds
  of kibibytes.
\end{example}

\subsection{Hierarchical register files}

The third and fourth observations suggest a hierarchy: values that are
read soon after being produced, and only once, can be served from a
small structure close to the ALUs.\tip{A value read once, by an
instruction issued soon after its producer, is the ideal cache
candidate: it is read and dies in the small file and never costs a
write to the main register file.}

\begin{definition}[Hierarchical register file]
  \label{def:l06-hierarchical-rf}
  A \term{hierarchical register file}[階層型レジスタファイル] places
  smaller, faster register storage in front of the
  \term{main register file (MRF)}[主レジスタファイル]\index{MRF|see{main register file (MRF)}}.  The two forms
  of \cref{fig:l06-rf-hierarchy} are a
  \term{register file cache (RFC)}[レジスタファイルキャッシュ]\index{RFC|see{register file cache (RFC)}}, one or
  two small levels (L0 and L1 RF) that hold recently produced values
  and filter reads and writes to the MRF, and a \emph{partitioned}
  register file with a fast and a slow segment, in which the slow,
  large segment can be built in a denser technology.
\end{definition}

\begin{figure}[tb]
  \centering
  % Two hierarchical register files (Lesson 6).  (a) A register file
% cache: small L0/L1 register files in front of the main register file,
% closest to the ALUs.  (b) A partitioned register file: a fast SRAM
% segment and a slow, large segment that may be built in non-volatile
% memory.
% Redraws the outline figure m6l2_opt1.
% Usage: % Two hierarchical register files (Lesson 6).  (a) A register file
% cache: small L0/L1 register files in front of the main register file,
% closest to the ALUs.  (b) A partitioned register file: a fast SRAM
% segment and a slow, large segment that may be built in non-volatile
% memory.
% Redraws the outline figure m6l2_opt1.
% Usage: % Two hierarchical register files (Lesson 6).  (a) A register file
% cache: small L0/L1 register files in front of the main register file,
% closest to the ALUs.  (b) A partitioned register file: a fast SRAM
% segment and a slow, large segment that may be built in non-volatile
% memory.
% Redraws the outline figure m6l2_opt1.
% Usage: \input{figures/l06-rf-hierarchy}
% Width: body (120 mm); the drawing is about 110 mm wide.
\begin{tikzpicture}[x=1mm, y=1mm, font=\scriptsize\sffamily,
    >={Stealth[length=1.6mm]}]
  \node[anchor=west, font=\scriptsize\sffamily\bfseries] at (0,40)
    {\iflangja{(a) レジスタファイルキャッシュ}{(a) register file cache}};
  \draw[fill=ln-box] (14,30) rectangle ++(20,6) node[midway] {\iflangja{演算器}{ALUs}};
  \draw[fill=ln-accent!35] (16,22) rectangle ++(16,6) node[midway] {L0 RF};
  \draw[fill=ln-accent!22] (10,14) rectangle ++(28,6) node[midway] {L1 RF};
  \draw[fill=ln-accent!10] (2,4) rectangle ++(44,8)
    node[midway] {\iflangja{主レジスタファイル (MRF)}{main register file (MRF)}};
  \draw[<->] (24,28) -- (24,30);
  \draw[<->] (24,20) -- (24,22);
  \draw[<->] (24,12) -- (24,14);
  \draw[->, ln-muted] (50,34) -- (50,6);
  \node[anchor=west, text=ln-muted] at (51,33) {\iflangja{小・高速}{small, fast}};
  \node[anchor=west, text=ln-muted] at (51,7) {\iflangja{大・低速}{large, slow}};
  % partitioned
  \node[anchor=west, font=\scriptsize\sffamily\bfseries] at (68,40)
    {\iflangja{(b) 分割型レジスタファイル}{(b) partitioned register file}};
  \draw[fill=ln-box] (76,30) rectangle ++(28,6) node[midway] {\iflangja{演算器}{ALUs}};
  \draw[fill=ln-accent!30] (70,14) rectangle ++(16,12)
    node[midway, align=center] {\iflangja{高速部\\(SRAM)}{fast\\segment\\(SRAM)}};
  \draw[fill=ln-tip!20] (86,14) rectangle ++(26,12)
    node[midway, align=center] {\iflangja{低速・大容量部\\(NVM も可)}{slow, large\\segment\\(may be NVM)}};
  \draw[<->] (82,26) -- (82,30);
  \draw[<->] (98,26) -- (98,30);
  \node[anchor=north, text=ln-muted, align=center] at (91,11)
    {\iflangja{よく使う値を高速部に置く}{frequently used values\\in the fast segment}};
\end{tikzpicture}

% Width: body (120 mm); the drawing is about 110 mm wide.
\begin{tikzpicture}[x=1mm, y=1mm, font=\scriptsize\sffamily,
    >={Stealth[length=1.6mm]}]
  \node[anchor=west, font=\scriptsize\sffamily\bfseries] at (0,40)
    {\iflangja{(a) レジスタファイルキャッシュ}{(a) register file cache}};
  \draw[fill=ln-box] (14,30) rectangle ++(20,6) node[midway] {\iflangja{演算器}{ALUs}};
  \draw[fill=ln-accent!35] (16,22) rectangle ++(16,6) node[midway] {L0 RF};
  \draw[fill=ln-accent!22] (10,14) rectangle ++(28,6) node[midway] {L1 RF};
  \draw[fill=ln-accent!10] (2,4) rectangle ++(44,8)
    node[midway] {\iflangja{主レジスタファイル (MRF)}{main register file (MRF)}};
  \draw[<->] (24,28) -- (24,30);
  \draw[<->] (24,20) -- (24,22);
  \draw[<->] (24,12) -- (24,14);
  \draw[->, ln-muted] (50,34) -- (50,6);
  \node[anchor=west, text=ln-muted] at (51,33) {\iflangja{小・高速}{small, fast}};
  \node[anchor=west, text=ln-muted] at (51,7) {\iflangja{大・低速}{large, slow}};
  % partitioned
  \node[anchor=west, font=\scriptsize\sffamily\bfseries] at (68,40)
    {\iflangja{(b) 分割型レジスタファイル}{(b) partitioned register file}};
  \draw[fill=ln-box] (76,30) rectangle ++(28,6) node[midway] {\iflangja{演算器}{ALUs}};
  \draw[fill=ln-accent!30] (70,14) rectangle ++(16,12)
    node[midway, align=center] {\iflangja{高速部\\(SRAM)}{fast\\segment\\(SRAM)}};
  \draw[fill=ln-tip!20] (86,14) rectangle ++(26,12)
    node[midway, align=center] {\iflangja{低速・大容量部\\(NVM も可)}{slow, large\\segment\\(may be NVM)}};
  \draw[<->] (82,26) -- (82,30);
  \draw[<->] (98,26) -- (98,30);
  \node[anchor=north, text=ln-muted, align=center] at (91,11)
    {\iflangja{よく使う値を高速部に置く}{frequently used values\\in the fast segment}};
\end{tikzpicture}

% Width: body (120 mm); the drawing is about 110 mm wide.
\begin{tikzpicture}[x=1mm, y=1mm, font=\scriptsize\sffamily,
    >={Stealth[length=1.6mm]}]
  \node[anchor=west, font=\scriptsize\sffamily\bfseries] at (0,40)
    {\iflangja{(a) レジスタファイルキャッシュ}{(a) register file cache}};
  \draw[fill=ln-box] (14,30) rectangle ++(20,6) node[midway] {\iflangja{演算器}{ALUs}};
  \draw[fill=ln-accent!35] (16,22) rectangle ++(16,6) node[midway] {L0 RF};
  \draw[fill=ln-accent!22] (10,14) rectangle ++(28,6) node[midway] {L1 RF};
  \draw[fill=ln-accent!10] (2,4) rectangle ++(44,8)
    node[midway] {\iflangja{主レジスタファイル (MRF)}{main register file (MRF)}};
  \draw[<->] (24,28) -- (24,30);
  \draw[<->] (24,20) -- (24,22);
  \draw[<->] (24,12) -- (24,14);
  \draw[->, ln-muted] (50,34) -- (50,6);
  \node[anchor=west, text=ln-muted] at (51,33) {\iflangja{小・高速}{small, fast}};
  \node[anchor=west, text=ln-muted] at (51,7) {\iflangja{大・低速}{large, slow}};
  % partitioned
  \node[anchor=west, font=\scriptsize\sffamily\bfseries] at (68,40)
    {\iflangja{(b) 分割型レジスタファイル}{(b) partitioned register file}};
  \draw[fill=ln-box] (76,30) rectangle ++(28,6) node[midway] {\iflangja{演算器}{ALUs}};
  \draw[fill=ln-accent!30] (70,14) rectangle ++(16,12)
    node[midway, align=center] {\iflangja{高速部\\(SRAM)}{fast\\segment\\(SRAM)}};
  \draw[fill=ln-tip!20] (86,14) rectangle ++(26,12)
    node[midway, align=center] {\iflangja{低速・大容量部\\(NVM も可)}{slow, large\\segment\\(may be NVM)}};
  \draw[<->] (82,26) -- (82,30);
  \draw[<->] (98,26) -- (98,30);
  \node[anchor=north, text=ln-muted, align=center] at (91,11)
    {\iflangja{よく使う値を高速部に置く}{frequently used values\\in the fast segment}};
\end{tikzpicture}

  \caption{Two hierarchical register files.  (a) A register file
    cache: small L0 and L1 register files between the ALUs and the main
    register file.  (b) A partitioned register file with a fast segment
    and a slow, large segment that may use non-volatile memory.}
  \label{fig:l06-rf-hierarchy}
\end{figure}

With an RFC, most reads hit in the small file, and a value that dies
before it is evicted never has to be written to the MRF.  The MRF can
then be smaller, or kept at the same size and accessed less often,
which saves energy because a small SRAM costs less per access than a
large one.\margin{Gebhart et al.: a 6-entry-per-thread RFC cuts MRF
reads by 50\,\% and writes by 59\,\%; with a two-level scheduler
that keeps a quarter of the threads active, register file energy falls
by 36\,\% \cite{gebhart2011energy}.}  GPUs make this easier than CPUs:
K12 §1.4.1 notes that GPUs do not support precise exceptions, so
\enquote{register files do not have to maintain architectural states}
\cite{kim2012gpgpu}, and a value can live only in the cache.

A candidate for the slow segment of a partitioned file is
\term{non-volatile memory (NVM)}[不揮発性メモリ]\index{NVM|see{non-volatile memory (NVM)}}, memory that keeps its
contents without power.  Built in NVM, the slow, large segment is
meant to save energy, while the values in frequent use stay in the
fast segment.\margin{NVM in a register
file is a research direction, not a shipped design; current GPUs build
the whole register file in SRAM.}

The caching decision can also be moved to the compiler, which knows the
lifetimes of values.  Gebhart, Keckler and Dally
\cite{gebhart2011compile} replace the hardware RFC by an
\term{operand register file (ORF)}[オペランドレジスタファイル]\index{ORF|see{operand register file (ORF)}}, a small
per-thread file whose contents the compiler allocates, and a
\term{last result file (LRF)}[ラストリザルトファイル]\index{LRF|see{last result file (LRF)}} of one entry per
thread that holds the most recent result for immediate reuse.\margin{The
compiler-managed hierarchy is reported to cut register file energy by up
to 54\,\%, against 34\,\% for hardware-only caching; in their baseline
the register file consumes 15--20\,\% of the processor's dynamic energy
\cite{gebhart2011compile}.}  Because the compiler places every value
explicitly, the ORF needs no tags and no hardware replacement, the same
property that distinguishes a scratchpad from a cache
(\cref{sec:l06-unified}).

\subsection{Register file virtualization}

The first two observations suggest sharing: if many registers are dead
or belong to finished threads at any moment, fewer physical registers
than names suffice \cite{jeon2015rfv}.

\begin{definition}[Architectural and physical registers]
  \label{def:l06-arch-phys}
  An \term{architectural register}[アーキテクチャレジスタ] is a
  register name as the program sees it (\texttt{R1} of thread~3); a
  \term{physical register}[物理レジスタ] is a storage entry of the
  register file.  The mapping of architectural to physical registers at
  run time is \term{register renaming}[レジスタリネーミング].  An
  out-of-order CPU
  renames a few architectural registers onto \emph{more} physical
  registers, to remove false dependences.
\end{definition}

\begin{definition}[Register file virtualization]
  \label{def:l06-rf-virtualization}
  The technique of
  \term{register file virtualization}[レジスタファイル仮想化] keeps the
  name space of all threads' architectural registers but maps it onto a
  \emph{smaller} physical register file: a physical register is assigned
  when a value is produced and released when the value dies, so that
  names whose values are live at different times share storage
  (\cref{fig:l06-rf-virtualization}).  The program semantics are
  unchanged; only the hardware is smaller.
\end{definition}

\begin{figure}[tb]
  \centering
  % Register file virtualisation (Lesson 6).  Six architectural register
% names (R1 of threads 1-3, R2 of threads 4-6) whose values are live in
% disjoint intervals share two physical registers: P0 and P1.
% Redraws the outline figure m6l2_opt2 (lifetime bars added).
% Usage: % Register file virtualisation (Lesson 6).  Six architectural register
% names (R1 of threads 1-3, R2 of threads 4-6) whose values are live in
% disjoint intervals share two physical registers: P0 and P1.
% Redraws the outline figure m6l2_opt2 (lifetime bars added).
% Usage: % Register file virtualisation (Lesson 6).  Six architectural register
% names (R1 of threads 1-3, R2 of threads 4-6) whose values are live in
% disjoint intervals share two physical registers: P0 and P1.
% Redraws the outline figure m6l2_opt2 (lifetime bars added).
% Usage: \input{figures/l06-rf-virtualization}
% Width: body (120 mm); the drawing is about 112 mm wide.
\begin{tikzpicture}[x=1mm, y=1mm, font=\scriptsize\sffamily,
    >={Stealth[length=1.6mm]}]
  \node[anchor=south west, font=\scriptsize\sffamily\bfseries] at (0,40)
    {\iflangja{名前空間}{name space}};
  \node[anchor=south west, font=\scriptsize\sffamily\bfseries] at (18,40)
    {\iflangja{値が生存する区間}{interval in which the value is live}};
  \node[anchor=south west, font=\scriptsize\sffamily\bfseries] at (90,40)
    {\iflangja{物理レジスタ}{physical registers}};
  \draw[ln-rule] (18,2) -- (82,2);
  \draw[->, ln-muted] (18,0) -- (82,0) node[right] {\iflangja{時間}{time}};
  \foreach \n/\a/\b/\c [count=\i from 0] in
      {Th1.R1/0/16/ln-accent!35, Th2.R1/20/36/ln-accent!35, Th3.R1/42/60/ln-accent!35,
       Th4.R2/6/22/ln-tip!35, Th5.R2/26/42/ln-tip!35, Th6.R2/46/62/ln-tip!35} {
    \node[anchor=west, font=\scriptsize\ttfamily] at (0,36-\i*6) {\n};
    \draw[ln-rule, dotted] (18,36-\i*6) -- (82,36-\i*6);
    \draw[fill=\c] (18+\a,34.7-\i*6) rectangle (18+\b,37.3-\i*6);
  }
  \draw[fill=ln-accent!35] (92,27) rectangle ++(14,6) node[midway] {P0};
  \draw[fill=ln-tip!35] (92,9) rectangle ++(14,6) node[midway] {P1};
  \draw[ln-muted] (83,38) -- ++(1.5,0) -- ++(0,-14) -- ++(-1.5,0);
  \draw[ln-muted] (83,20) -- ++(1.5,0) -- ++(0,-14) -- ++(-1.5,0);
  \draw[->] (84.5,31) -- (92,30);
  \draw[->] (84.5,13) -- (92,12);
\end{tikzpicture}

% Width: body (120 mm); the drawing is about 112 mm wide.
\begin{tikzpicture}[x=1mm, y=1mm, font=\scriptsize\sffamily,
    >={Stealth[length=1.6mm]}]
  \node[anchor=south west, font=\scriptsize\sffamily\bfseries] at (0,40)
    {\iflangja{名前空間}{name space}};
  \node[anchor=south west, font=\scriptsize\sffamily\bfseries] at (18,40)
    {\iflangja{値が生存する区間}{interval in which the value is live}};
  \node[anchor=south west, font=\scriptsize\sffamily\bfseries] at (90,40)
    {\iflangja{物理レジスタ}{physical registers}};
  \draw[ln-rule] (18,2) -- (82,2);
  \draw[->, ln-muted] (18,0) -- (82,0) node[right] {\iflangja{時間}{time}};
  \foreach \n/\a/\b/\c [count=\i from 0] in
      {Th1.R1/0/16/ln-accent!35, Th2.R1/20/36/ln-accent!35, Th3.R1/42/60/ln-accent!35,
       Th4.R2/6/22/ln-tip!35, Th5.R2/26/42/ln-tip!35, Th6.R2/46/62/ln-tip!35} {
    \node[anchor=west, font=\scriptsize\ttfamily] at (0,36-\i*6) {\n};
    \draw[ln-rule, dotted] (18,36-\i*6) -- (82,36-\i*6);
    \draw[fill=\c] (18+\a,34.7-\i*6) rectangle (18+\b,37.3-\i*6);
  }
  \draw[fill=ln-accent!35] (92,27) rectangle ++(14,6) node[midway] {P0};
  \draw[fill=ln-tip!35] (92,9) rectangle ++(14,6) node[midway] {P1};
  \draw[ln-muted] (83,38) -- ++(1.5,0) -- ++(0,-14) -- ++(-1.5,0);
  \draw[ln-muted] (83,20) -- ++(1.5,0) -- ++(0,-14) -- ++(-1.5,0);
  \draw[->] (84.5,31) -- (92,30);
  \draw[->] (84.5,13) -- (92,12);
\end{tikzpicture}

% Width: body (120 mm); the drawing is about 112 mm wide.
\begin{tikzpicture}[x=1mm, y=1mm, font=\scriptsize\sffamily,
    >={Stealth[length=1.6mm]}]
  \node[anchor=south west, font=\scriptsize\sffamily\bfseries] at (0,40)
    {\iflangja{名前空間}{name space}};
  \node[anchor=south west, font=\scriptsize\sffamily\bfseries] at (18,40)
    {\iflangja{値が生存する区間}{interval in which the value is live}};
  \node[anchor=south west, font=\scriptsize\sffamily\bfseries] at (90,40)
    {\iflangja{物理レジスタ}{physical registers}};
  \draw[ln-rule] (18,2) -- (82,2);
  \draw[->, ln-muted] (18,0) -- (82,0) node[right] {\iflangja{時間}{time}};
  \foreach \n/\a/\b/\c [count=\i from 0] in
      {Th1.R1/0/16/ln-accent!35, Th2.R1/20/36/ln-accent!35, Th3.R1/42/60/ln-accent!35,
       Th4.R2/6/22/ln-tip!35, Th5.R2/26/42/ln-tip!35, Th6.R2/46/62/ln-tip!35} {
    \node[anchor=west, font=\scriptsize\ttfamily] at (0,36-\i*6) {\n};
    \draw[ln-rule, dotted] (18,36-\i*6) -- (82,36-\i*6);
    \draw[fill=\c] (18+\a,34.7-\i*6) rectangle (18+\b,37.3-\i*6);
  }
  \draw[fill=ln-accent!35] (92,27) rectangle ++(14,6) node[midway] {P0};
  \draw[fill=ln-tip!35] (92,9) rectangle ++(14,6) node[midway] {P1};
  \draw[ln-muted] (83,38) -- ++(1.5,0) -- ++(0,-14) -- ++(-1.5,0);
  \draw[ln-muted] (83,20) -- ++(1.5,0) -- ++(0,-14) -- ++(-1.5,0);
  \draw[->] (84.5,31) -- (92,30);
  \draw[->] (84.5,13) -- (92,12);
\end{tikzpicture}

  \caption{Register file virtualization.  The values of \texttt{R1} in
    threads 1--3 are live in disjoint intervals and share physical
    register P0; those of \texttt{R2} in threads 4--6 share P1.  Six
    names need two physical registers.}
  \label{fig:l06-rf-virtualization}
\end{figure}

\caution{The direction is the opposite of CPU renaming: a CPU has more
physical than architectural registers, register file virtualization on
a GPU has fewer.  A virtualized file is resource sharing: only the
name space has to be maintained.}%
The mapping costs a table lookup per operand, and the design must
release a register exactly when its value dies, which requires lifetime
information from the compiler (the last read of each value).  It must
also guarantee progress when the physical registers run out: a warp
that needs a new register while all are taken must wait for others to
release theirs, and if every warp waits, none progresses.

\begin{example}
  \label{ex:l06-rf-virtualization}
  A kernel uses 32 registers per thread and the SM keeps 1024 threads
  resident, so the name space holds $32 \times 1024 = \num{32768}$
  registers, \SI{128}{\kibi\byte} at \SI{4}{\byte} each.  Suppose
  profiling shows that at most 60\,\% of them hold live values at any
  moment, partly because of dead registers and partly because of
  finished threads.  A virtualized file of
  $0.6 \times \num{32768} \approx \num{19661}$ physical registers,
  about \SI{77}{\kibi\byte}, then holds every live value; the remaining
  \SI{51}{\kibi\byte} can be removed, or kept to raise the number of
  resident threads.\margin{Our numbers, chosen for the arithmetic; no
  measured live fraction is implied.}
\end{example}

\begin{keypoint}
  The GPU register file is large because it holds every resident
  thread's registers, and must be fast because it feeds every lane.
  Most values are dead for long intervals or read once soon after
  being produced, which justifies a hierarchy (register file cache,
  ORF/LRF, fast and slow partitions) for latency and energy, and
  virtualization (fewer physical registers than names) for capacity.
\end{keypoint}

\section{Unified register file and shared memory}
\label{sec:l06-unified}

Lesson~3 showed that occupancy is limited by whichever resource a
kernel exhausts first: thread slots, registers or shared
memory.\source{M6, unified register file and shared memory;
\cite{gebhart2012unifying}; A18 ch.~4; CPG Compute Capability 7.x--9.0,
Shared Memory.}  Kernels differ widely in how they use the latter two,
so a fixed split between them wastes storage.  This section compares
the three kinds of on-chip storage and presents the unified design.

\subsection{Occupancy revisited}

Consider an SM with a \SI{128}{\kibi\byte} register file (\num{32768}
registers) and \SI{64}{\kibi\byte} of shared memory,\margin{A
configurable split predates the unified design: a Fermi SM (2009) splits
\SI{64}{\kibi\byte} into \SI{48}{\kibi\byte} of shared memory and
\SI{16}{\kibi\byte} of L1 or the reverse \cite{nvidia2009fermi}.} and two kernels
launched with 256 threads per block (\cref{fig:l06-unified-pool}).
\begin{itemize}
  \item Kernel~A uses 64 registers per thread and no shared memory:
    $64 \times 256 \times \SI{4}{\byte} = \SI{64}{\kibi\byte}$ of
    registers per block, so two blocks fill the register file and the
    whole shared memory idles.
  \item Kernel~B uses 16 registers per thread and \SI{16}{\kibi\byte}
    of shared memory per block: \SI{16}{\kibi\byte} of registers per
    block.  Shared memory admits four blocks, which use only half of the
    register file.\margin{The limit
on resident threads (2048 per SM on compute capabilities 7.0, 8.0 and
9.0, CPG Table~21 \cite{nvidia2024cudaguide}) still applies; kernel~B's
six blocks are 1536 threads.}
\end{itemize}

\begin{figure}[tb]
  \centering
  % Split versus unified on-chip storage (Lesson 6, own example).  An SM
% with a 128 KiB register file and 64 KiB of shared memory, or one
% 192 KiB pool.  Kernel A: 64 KiB of registers per block, no shared
% memory.  Kernel B: 16 KiB of registers and 16 KiB of shared memory per
% block.  Each filled segment is one resident block's allocation; white
% is idle.  Scale 0.45 mm per KiB.
% Original (no outline figure for this lesson).
% Usage: % Split versus unified on-chip storage (Lesson 6, own example).  An SM
% with a 128 KiB register file and 64 KiB of shared memory, or one
% 192 KiB pool.  Kernel A: 64 KiB of registers per block, no shared
% memory.  Kernel B: 16 KiB of registers and 16 KiB of shared memory per
% block.  Each filled segment is one resident block's allocation; white
% is idle.  Scale 0.45 mm per KiB.
% Original (no outline figure for this lesson).
% Usage: % Split versus unified on-chip storage (Lesson 6, own example).  An SM
% with a 128 KiB register file and 64 KiB of shared memory, or one
% 192 KiB pool.  Kernel A: 64 KiB of registers per block, no shared
% memory.  Kernel B: 16 KiB of registers and 16 KiB of shared memory per
% block.  Each filled segment is one resident block's allocation; white
% is idle.  Scale 0.45 mm per KiB.
% Original (no outline figure for this lesson).
% Usage: \input{figures/l06-unified-pool}
% Width: body (120 mm); the drawing is about 116 mm wide.
\begin{tikzpicture}[x=1mm, y=1mm, font=\scriptsize\sffamily]
  \def\k{0.45}
  % \seg{x0 KiB}{y}{size KiB}{fill}
  \newcommand{\seg}[4]{\draw[fill=#4] ({26+(#1)*\k},#2) rectangle ++({(#3)*\k},4);}
  \newcommand{\idle}[3]{\draw[ln-rule, fill=white] ({26+(#1)*\k},#2) rectangle ++({(#3)*\k},4);}
  % headings
  \node[anchor=south west, text=ln-muted] at (26,46) {\iflangja{レジスタファイル 128\,KiB}{register file 128 KiB}};
  \node[anchor=south west, text=ln-muted] at (26+131*\k,46) {\iflangja{共有メモリ 64\,KiB}{shared 64 KiB}};
  % Kernel A split: 2 blocks x 64 KiB regs
  \node[anchor=west] at (0,42) {\iflangja{A：分離}{A, split}};
  \seg{0}{40}{64}{ln-accent!30} \seg{64}{40}{64}{ln-accent!30}
  \idle{131}{40}{64}
  \node[anchor=west] at (26+196*\k,42) {2};
  % Kernel A unified
  \node[anchor=west] at (0,34) {\iflangja{A：統合}{A, unified}};
  \seg{0}{32}{64}{ln-accent!30} \seg{64}{32}{64}{ln-accent!30} \seg{128}{32}{64}{ln-accent!30}
  \node[anchor=west] at (26+196*\k,34) {3};
  % Kernel B split: 4 blocks x 16 regs + 4 x 16 smem
  \node[anchor=west] at (0,22) {\iflangja{B：分離}{B, split}};
  \foreach \b in {0,...,3} { \seg{\b*16}{20}{16}{ln-accent!30} \seg{131+\b*16}{20}{16}{ln-tip!35} }
  \idle{64}{20}{64}
  \node[anchor=west] at (26+196*\k,22) {4};
  % Kernel B unified: 6 blocks x (16 + 16)
  \node[anchor=west] at (0,14) {\iflangja{B：統合}{B, unified}};
  \foreach \b in {0,...,5} { \seg{\b*32}{12}{16}{ln-accent!30} \seg{\b*32+16}{12}{16}{ln-tip!35} }
  \node[anchor=west] at (26+196*\k,14) {6};
  \node[anchor=south, text=ln-muted] at (26+198*\k,46) {\iflangja{ブロック}{blocks}};
  % legend
  \draw[fill=ln-accent!30] (26,3) rectangle ++(4,3);
  \node[anchor=west] at (31,4.5) {\iflangja{レジスタ}{registers}};
  \draw[fill=ln-tip!35] (50,3) rectangle ++(4,3);
  \node[anchor=west] at (55,4.5) {\iflangja{共有メモリ}{shared memory}};
  \draw[ln-rule, fill=white] (78,3) rectangle ++(4,3);
  \node[anchor=west] at (83,4.5) {\iflangja{未使用}{idle}};
  \node[anchor=south west, text=ln-muted] at (26,26.5)
    {\iflangja{統合プール 192\,KiB（1 本の帯）}{unified pool 192 KiB (one bar)}};
\end{tikzpicture}

% Width: body (120 mm); the drawing is about 116 mm wide.
\begin{tikzpicture}[x=1mm, y=1mm, font=\scriptsize\sffamily]
  \def\k{0.45}
  % \seg{x0 KiB}{y}{size KiB}{fill}
  \newcommand{\seg}[4]{\draw[fill=#4] ({26+(#1)*\k},#2) rectangle ++({(#3)*\k},4);}
  \newcommand{\idle}[3]{\draw[ln-rule, fill=white] ({26+(#1)*\k},#2) rectangle ++({(#3)*\k},4);}
  % headings
  \node[anchor=south west, text=ln-muted] at (26,46) {\iflangja{レジスタファイル 128\,KiB}{register file 128 KiB}};
  \node[anchor=south west, text=ln-muted] at (26+131*\k,46) {\iflangja{共有メモリ 64\,KiB}{shared 64 KiB}};
  % Kernel A split: 2 blocks x 64 KiB regs
  \node[anchor=west] at (0,42) {\iflangja{A：分離}{A, split}};
  \seg{0}{40}{64}{ln-accent!30} \seg{64}{40}{64}{ln-accent!30}
  \idle{131}{40}{64}
  \node[anchor=west] at (26+196*\k,42) {2};
  % Kernel A unified
  \node[anchor=west] at (0,34) {\iflangja{A：統合}{A, unified}};
  \seg{0}{32}{64}{ln-accent!30} \seg{64}{32}{64}{ln-accent!30} \seg{128}{32}{64}{ln-accent!30}
  \node[anchor=west] at (26+196*\k,34) {3};
  % Kernel B split: 4 blocks x 16 regs + 4 x 16 smem
  \node[anchor=west] at (0,22) {\iflangja{B：分離}{B, split}};
  \foreach \b in {0,...,3} { \seg{\b*16}{20}{16}{ln-accent!30} \seg{131+\b*16}{20}{16}{ln-tip!35} }
  \idle{64}{20}{64}
  \node[anchor=west] at (26+196*\k,22) {4};
  % Kernel B unified: 6 blocks x (16 + 16)
  \node[anchor=west] at (0,14) {\iflangja{B：統合}{B, unified}};
  \foreach \b in {0,...,5} { \seg{\b*32}{12}{16}{ln-accent!30} \seg{\b*32+16}{12}{16}{ln-tip!35} }
  \node[anchor=west] at (26+196*\k,14) {6};
  \node[anchor=south, text=ln-muted] at (26+198*\k,46) {\iflangja{ブロック}{blocks}};
  % legend
  \draw[fill=ln-accent!30] (26,3) rectangle ++(4,3);
  \node[anchor=west] at (31,4.5) {\iflangja{レジスタ}{registers}};
  \draw[fill=ln-tip!35] (50,3) rectangle ++(4,3);
  \node[anchor=west] at (55,4.5) {\iflangja{共有メモリ}{shared memory}};
  \draw[ln-rule, fill=white] (78,3) rectangle ++(4,3);
  \node[anchor=west] at (83,4.5) {\iflangja{未使用}{idle}};
  \node[anchor=south west, text=ln-muted] at (26,26.5)
    {\iflangja{統合プール 192\,KiB（1 本の帯）}{unified pool 192 KiB (one bar)}};
\end{tikzpicture}

% Width: body (120 mm); the drawing is about 116 mm wide.
\begin{tikzpicture}[x=1mm, y=1mm, font=\scriptsize\sffamily]
  \def\k{0.45}
  % \seg{x0 KiB}{y}{size KiB}{fill}
  \newcommand{\seg}[4]{\draw[fill=#4] ({26+(#1)*\k},#2) rectangle ++({(#3)*\k},4);}
  \newcommand{\idle}[3]{\draw[ln-rule, fill=white] ({26+(#1)*\k},#2) rectangle ++({(#3)*\k},4);}
  % headings
  \node[anchor=south west, text=ln-muted] at (26,46) {\iflangja{レジスタファイル 128\,KiB}{register file 128 KiB}};
  \node[anchor=south west, text=ln-muted] at (26+131*\k,46) {\iflangja{共有メモリ 64\,KiB}{shared 64 KiB}};
  % Kernel A split: 2 blocks x 64 KiB regs
  \node[anchor=west] at (0,42) {\iflangja{A：分離}{A, split}};
  \seg{0}{40}{64}{ln-accent!30} \seg{64}{40}{64}{ln-accent!30}
  \idle{131}{40}{64}
  \node[anchor=west] at (26+196*\k,42) {2};
  % Kernel A unified
  \node[anchor=west] at (0,34) {\iflangja{A：統合}{A, unified}};
  \seg{0}{32}{64}{ln-accent!30} \seg{64}{32}{64}{ln-accent!30} \seg{128}{32}{64}{ln-accent!30}
  \node[anchor=west] at (26+196*\k,34) {3};
  % Kernel B split: 4 blocks x 16 regs + 4 x 16 smem
  \node[anchor=west] at (0,22) {\iflangja{B：分離}{B, split}};
  \foreach \b in {0,...,3} { \seg{\b*16}{20}{16}{ln-accent!30} \seg{131+\b*16}{20}{16}{ln-tip!35} }
  \idle{64}{20}{64}
  \node[anchor=west] at (26+196*\k,22) {4};
  % Kernel B unified: 6 blocks x (16 + 16)
  \node[anchor=west] at (0,14) {\iflangja{B：統合}{B, unified}};
  \foreach \b in {0,...,5} { \seg{\b*32}{12}{16}{ln-accent!30} \seg{\b*32+16}{12}{16}{ln-tip!35} }
  \node[anchor=west] at (26+196*\k,14) {6};
  \node[anchor=south, text=ln-muted] at (26+198*\k,46) {\iflangja{ブロック}{blocks}};
  % legend
  \draw[fill=ln-accent!30] (26,3) rectangle ++(4,3);
  \node[anchor=west] at (31,4.5) {\iflangja{レジスタ}{registers}};
  \draw[fill=ln-tip!35] (50,3) rectangle ++(4,3);
  \node[anchor=west] at (55,4.5) {\iflangja{共有メモリ}{shared memory}};
  \draw[ln-rule, fill=white] (78,3) rectangle ++(4,3);
  \node[anchor=west] at (83,4.5) {\iflangja{未使用}{idle}};
  \node[anchor=south west, text=ln-muted] at (26,26.5)
    {\iflangja{統合プール 192\,KiB（1 本の帯）}{unified pool 192 KiB (one bar)}};
\end{tikzpicture}

  \caption{Split versus unified storage for two kernels with 256
    threads per block.  Kernel~A fills the register file and leaves
    shared memory idle; kernel~B fills shared memory and leaves half
    the register file idle.  One \SI{192}{\kibi\byte} pool admits three
    blocks of A and six of B.}
  \label{fig:l06-unified-pool}
\end{figure}

If the \SI{192}{\kibi\byte} were one pool from which each block takes
its registers and its shared memory, kernel~A would fit
$192/64 = 3$ blocks and kernel~B $192/32 = 6$ blocks: 50\,\% more
resident threads for both (\cref{tab:l06-unified}).  In the unified design, occupancy is
limited by registers and shared memory \emph{together}.

\begin{margintable}
  \caption{Resident blocks (threads) per SM for the kernels of
    \cref{fig:l06-unified-pool}.}
  \label{tab:l06-unified}
  \begin{tabular}{@{}lcc@{}}
    \toprule
    kernel & split & unified \\
    \midrule
    A & 2 (512) & 3 (768) \\
    B & 4 (1024) & 6 (1536) \\
    \bottomrule
  \end{tabular}
\end{margintable}

\subsection{Caches, scratchpads and register files}

Whether registers and shared memory can share storage depends on how
each is addressed.  A cache and a scratchpad differ in who decides what
they hold.

\begin{definition}[Hardware- and software-managed caches]
  \label{def:l06-managed-caches}
  A \term{hardware-managed cache}[ハードウェア管理キャッシュ] decides by
  itself which memory blocks it holds; each entry therefore carries a
  \term{tag}[タグ], the upper address bits of the block it holds,
  compared on every access to tell whether the requested address is
  present, and a full cache
  evicts a block chosen by its
  \term{replacement policy}[置換ポリシ], for example
  \term{LRU (least recently used)}[LRU]\index{least recently used|see{LRU (least recently used)}}, which evicts the block whose
  last access lies furthest in the past.  A
  \term{software-managed cache}[ソフトウェア管理キャッシュ] (the
  scratchpad memory of Lesson~3, CUDA's shared memory) holds what the
  program puts there: the program knows where every datum lives, so
  there are no tags, no misses and no replacement.
\end{definition}
\margin{NVIDIA GPUs cache global memory in \SI{128}{\byte} lines of
four \SI{32}{\byte} sectors (Lesson~4), and a warp's accesses are
served in \SI{32}{\byte} transactions (BPG §10.2.1 Coalesced Access to
Global Memory \cite{nvidia2026cudabestpractices}).}

A register file resembles a scratchpad: both store data only, have no
tags, and select an entry by its location rather than by a memory
address.  They differ in two respects (\cref{tab:l06-storage}).\margin{Shared memory latency is roughly 100 times lower than that of
uncached global memory \cite{harris2013sharedmem}.}  An
arithmetic instruction reads two source registers, three for a fused
multiply-add $d = a \cdot b + c$, and writes one, so a
register file needs more ports than a memory, which serves one address
per access.  And the registers an instruction uses are named in the
instruction, so its access pattern is known statically, before
execution, whereas a scratchpad address is computed at run time.  The
static knowledge is what allows the bank layouts of Lesson~4 and the
compiler-managed ORF above, despite the higher bandwidth.

\begin{table}[tb]
  \centering
  \caption{Three kinds of on-chip storage.}
  \label{tab:l06-storage}
  \setlength{\tabcolsep}{4pt}
  \begin{tabular}{@{}>{\raggedright\arraybackslash}p{30mm}
                  >{\raggedright\arraybackslash}p{22mm}
                  >{\raggedright\arraybackslash}p{24mm}
                  >{\raggedright\arraybackslash}p{36mm}@{}}
    \toprule
    & hardware cache (L1) & scratchpad (shared memory) & register file \\
    \midrule
    contents chosen by & hardware & program & compiler \\
    tags & yes & no & no \\
    located by & address and tag & location & location (register name) \\
    \addlinespace
    accesses per instruction & one address & one address & 2--3 reads, one write \\
    pattern known & at run time & at run time & at compile time \\
    \bottomrule
  \end{tabular}
\end{table}

\subsection{A unified structure}

Since all three structures are SRAM arrays on the same SM, one array
can serve all of them, partitioned per kernel.\margin{Gebhart et al.\ build on
their earlier register hierarchy and two-level scheduler, which
\enquote{reduce accesses to the main register file by 60\,\%}, so that
a pool with fewer ports per byte than a dedicated register file can
still supply the operands \cite{gebhart2012unifying}.}

\begin{definition}[Unified local memory]
  \label{def:l06-unified-local-memory}
  The design called \term{unified local memory}[統合ローカルメモリ] is a
  single on-chip
  SRAM pool per SM whose capacity is divided, for each kernel, among the
  register file, the shared memory and, in the full proposal, the
  \term{L1 data cache}[L1データキャッシュ], the hardware-managed
  first-level cache for global and local memory accesses.
\end{definition}

The unified pool must meet the strictest demand of each user: the
bandwidth of the register file, hence enough read and write ports, and
more banks than either structure alone, because register accesses (the
same register in every lane) and shared-memory accesses (by computed
address) have different patterns and must both find their banks free.
In exchange, each kernel gets the split it needs, relaxing its
occupancy limit.

\caution{Lecture vs paper: this chapter, following the lecture,
unifies the register file and shared memory.  Gebhart et al.\ unify \emph{three} structures,
\enquote{the register file, scratchpad (shared memory in NVIDIA
terminology), and cache}, and report \enquote{a performance benefit as
high as 71\,\% along with an energy reduction up to 33\,\%}
\cite{gebhart2012unifying}.}%
Shipping GPUs share one array only between the L1 data cache and
shared memory, as Fermi did and Volta and later GPUs do again, while
the register file stays separate.

\begin{definition}[Shared memory carve-out]
  \label{def:l06-carveout}
  The \term{shared memory carve-out}[共有メモリのカーブアウト] is the
  part of an SM's combined L1 data cache and shared memory array that
  is configured as shared memory; the rest serves as L1 cache.  The
  CUDA runtime chooses it per kernel from a fixed set of sizes, guided
  by the kernel's shared memory use and an optional preference.
\end{definition}
\tip{A kernel that needs more
than \SI{48}{\kibi\byte} of dynamic shared memory per block must opt in
with \texttt{cudaFuncAttribute\allowbreak MaxDynamic\allowbreak SharedMemorySize} (CPG,
Compute Capability 7.x, Shared Memory \cite{nvidia2024cudaguide}).}%
The preference is set per kernel with one runtime call.

\begin{lstlisting}[language=C++, numbers=none]
// Prefer about half of the maximum shared memory
// capacity for kernel k (a hint; the runtime rounds
// up to a supported carve-out).
cudaFuncSetAttribute(k,
    cudaFuncAttributePreferredSharedMemoryCarveout, 50);
\end{lstlisting}

\margin{Combined L1 and shared memory per SM: \SI{128}{\kibi\byte} on
V100, \SI{192}{\kibi\byte} on A100, \SI{256}{\kibi\byte} on H100.
Carve-outs: V100 0, 8, 16, 32, 64 or \SI{96}{\kibi\byte}; A100 up to
\SI{164}{\kibi\byte}; H100 0, 8, 16, 32, 64, 100, 132, 164, 196 or
\SI{228}{\kibi\byte} (CPG, Compute Capability 7.x, 8.x and 9.0
\cite{nvidia2024cudaguide}).}%
A kernel that uses little shared memory thus gets a large L1, one that
tiles its data gets most of the array as scratchpad: the flexibility of
\cref{fig:l06-unified-pool} for two structures.

\begin{keypoint}[Lesson summary]
  A divergent branch makes a warp run its paths one after the other
  under partial masks until they reconverge at the immediate
  post-dominator.  If-conversion turns the branch into predicated code
  that fetches both paths; the SIMT stack pushes (alternate PC, mask,
  reconvergence PC) entries at divergent, also nested, branches and pops
  them at the reconvergence PC.  Wider warps save
  fetches but diverge more; dynamic warp formation merges threads of
  different warps at one PC, each in its home lane.  Dead, short-lived
  and read-once register values justify register file caches, ORF/LRF
  levels and fast/slow partitions; virtualization maps the names onto
  fewer physical registers.
  Unified local memory partitions one pool among registers, shared
  memory and the L1 per kernel; current GPUs do so for L1 and shared
  memory through the carve-out.
\end{keypoint}

\end{document}
